\section{AIME I 2017}
        \begin{enumerate}
        	\item (\textit{AIME I 2017}) Fifteen distinct points are designated on $\triangle ABC$: the 3 vertices $A$, $B$, and $C$; $3$ other points on side $\overline{AB}$; $4$ other points on side $\overline{BC}$; and $5$ other points on side $\overline{CA}$. Find the number of triangles with positive area whose vertices are among these $15$ points.


 
 \textbf{Answer: }390



\textbf{Solution 1}

Every triangle is uniquely determined by 3 points. There are $\binom{15}{3}=455$ ways to choose 3 points, but that counts the degenerate triangles formed by choosing three points on a line. There are $\binom{5}{3}$ invalid cases on segment $AB$, $\binom{6}{3}$ invalid cases on segment $BC$, and $\binom{7}{3}$ invalid cases on segment $CA$ for a total of $65$ invalid cases. The answer is thus $455-65=\boxed{390}$.

 



\textbf{Solution 2}

We simply choose $3$ points from $15$ to begin with, since a triangle consists of $3$ points and there are $15$ points total. That gives us $\binom{15}{3}$. Now, we need to remove the degenerate triangles from this total. The degenerate triangles are just the triangles that are collinear. We can count the number of degenerate triangles on each side of $ABC$, then subtract it from $\binom{15}{3}$. For $\overline{AB}$, we have $\binom{5}{3}$ degenerate triangles (don't forget $A$ and $B$). Then, for $\overline{BC}$, we have $\binom{6}{3}$ degenerate triangles, and for $\overline{AC}$, we have $\binom{7}{3}$ degenerate triangles. So, our final answer is $\binom{15}{3} - \binom{7}{3} - \binom{6}{3} - \binom{5}{3} = \boxed{390}$

 -jb2015007 and cnwyshi (minor edits like grammar)

 



\textbf{Solution 3 (Casework Ba\$h)}

We will use casework based on the number of common vertices to solve this problem.

 Case 1: No common vertices

 There are two subcases for this solution, being three distinct sides, and two vertices on a side. This is because, for a triangle to have area 0, all three points must be on one side, so this scenario is possible.

 First, to calculate the amount of triangles where all three vertices are on 3 distinct sides of a triangle, one just uses multiplication, for $5\times4\times3=60$.

 The other case is a bit trickier, as there are three possible sides where the two common vertices could be; AC, AB and BC.

 If the two vertices are on segment AC, where there are 5 points to choose from, the last point must be on a different segment. These can be combined in any way, so there are $(5, 2) \times 7$ ways to choose this case, or $70$.

 If the two vertices are on segment AB, there are 3 points to choose two from, and 9 other points. This yields $(3, 2) \times 9 = 3 \times 9 = 27$ possibilities.

 Finally, if the two vertices are on BC, there are 4 points to choose two from, and 8 others. This yields $(4, 2) \times 8 = 6 \times 8 = 48$ possibilities.

 Adding all of the zero vertice cases together results in $60 + 70 + 27 + 48 = 205$ non-common-vertice solutions.

 Case 2: One Common Vertice

 There are three cases to work with here; one where the vertice is A, one where it is B, and one where it is C.

 If the common vertice is A, and all other vertices are on distinct line segments (AB, BC, or CA), there is guaranteed to be a line segment where a vertice does not occur. This means that there are three subcases for this type, relative to the large case of a vertice on A: a vertice not being on AB (meaning vertices on BC and AC), a vertice not being on BC (meaning vertices on AC and AB) and a vertice not on AC (meaning vertices on AB and BC).

 For these subcases, only multiplying the number of points on each line suffices, so for this case for A specifically, we have $(3\times5)+(4\times5)+(3\times4)=15+20+12=47$ ways.

 This is true for all one-vertice distinct line segment cases (because of symmetry), so this case can be replicated three times, for $47 \times 3 = 141$ cases of 1-vertice distinct line segment.

 However, it is also possible that two points are on the same line segment, which must be directly opposite to the vertice (because no 3 points can be on the same line). This means that there must be 2 dots on that line segment opposite the vertice. The line segments opposite vertice A, B, and C are BC, AC and AB respectively, which have 4, 5, and 3 dots respectively. This problem asks for us to choose 2 from four, five, and 3 and add them up, which gets $(4,2)+(5,2)+(3,2)=6+10+3=19$ cases.

 In total, for one vertice cases, we have $141+19=160$ cases.This brings the aggregate number of cases to $205+160=365$ cases.

 Case 3: Two Common Vertices

 The only way a triangle could not have an area of greater than 0 is if all three of its vertices are on the same line. Two are already on the same line, so the last point cannot be on the same line as the other two. This is also the two vertice case, so the final point cannot be a vertice. Out of the 12 non-vertice points, 3, 4, or 5 of them will be unavailable, whether if B, A, or C is chosen, respectively. This gives $9+8+7=24$ cases for two-vertice solutions.

 Adding this to the aggregate gives $365+24=389$ cases.

 Case 4: Three Common Vertices

 By trivial, this is the original triangle, which is still a case. This gives 1 three-vertice solution.

 This brings the total number to $390$ cases, which is the final solution. $\boxed{390}$.

 


	\item (\textit{AIME I 2017}) When each of $702$, $787$, and $855$ is divided by the positive integer $m$, the remainder is always the positive integer $r$. When each of $412$, $722$, and $815$ is divided by the positive integer $n$, the remainder is always the positive integer $s \neq r$. Find $m+n+r+s$.


 
 \textbf{Answer: }062



\textbf{Solution}

Let's work on both parts of the problem separately. First, 
 \[855 \equiv 787 \equiv 702 \equiv r \pmod{m}.\] We take the difference of $855$ and $787$, and also of $787$ and $702$. We find that they are $85$ and $68$, respectively. Since the greatest common divisor of the two differences is $17$, $m=1,17$ but if $m=1$, then $r=0$ which violates $r$ being positive,so it's safe to assume that $m = 17$.

 Then, we divide $855$ by $17$, and it's easy to see that $r = 5$. Dividing $787$ and $702$ by $17$ also yields remainders of $5$, which means our work up to here is correct.

 Doing the same thing with $815$, $722$, and $412$, the differences between $815$ and $722$ and $412$ are $310$ and $93$, respectively. Since the only common divisor (besides $1$, of course) is $31$, $n = 31$. Dividing all $3$ numbers by $31$ yields a remainder of $9$ for each, so $s = 9$. Thus, $m + n + r + s = 17 + 5 + 31 + 9 = \boxed{062}$.

 



\textbf{Solution 2}

We know that $702 = am + r, 787 = bm + r,$ and $855 = cm+r$ where $a-c$ are integers.

 Subtracting the first two, the first and third, and the last two, we get $85 = (b-a)m, 153=(c-a)m,$ and $68=(c-b)m.$

 We know that $b-a, c-a$ and $c-b$ must be integers, so all the numbers are divisible by $m.$ 

 Factorizing the numbers, we get $85 = 5 \cdot 17, 153 = 3^2 \cdot 17,$ and $68 = 2^2 \cdot 17.$ We see that all these have a factor of 17, so $m=17.$

 Finding the remainder when $702$ is divided by $17,$ we get $n=5.$

 Doing the same thing for the next three numbers, we get $17 + 5 + 31 + 9 = \boxed{062}$




\textbf{Solution 3 (Sol. 1 but possibly more clear)}

As in Solution 1, we are given $855\equiv787\equiv702\equiv r\pmod{m}$ and $815\equiv722\equiv412\equiv s\pmod{n}.$ Tackling the first equation, we can simply look at $855\equiv787\equiv702\pmod{m}$. We subtract $702$ from each component of the congruency to get $153\equiv85\equiv0\pmod{m}$. Thus, we know that $153$ and $85$ must both be divisible by $m$. The only possible $m$, in this case, become $17$ and $1$; obviously, $m\neq1$, so we know $m=17$. We go back to the original equation, plug in $m$, and we find that $r=5$. 

 Similarly, we can subtract out the smallest value in the second congruency, $412$. We end up with $403\equiv310\equiv0\pmod{n}$. Again, we find that $n=31$ or $1$, so $n=31$. We also find that $s=9$. 

 Thus, our answer is $\boxed{062}$. 

 


	\item (\textit{AIME I 2017}) For a positive integer $n$, let $d_n$ be the units digit of $1 + 2 + \dots + n$. Find the remainder when

 \[\sum_{n=1}^{2017} d_n\]is divided by $1000$.


 
 \textbf{Answer: }069



\textbf{Solution}

We see that $d_n$ appears in cycles of $20$ and the cycles are 
 \[1,3,6,0,5,1,8,6,5,5,6,8,1,5,0,6,3,1,0,0,\] adding a total of $70$ each cycle.Since $\left\lfloor\frac{2017}{20}\right\rfloor=100$, we know that by $2017$, there have been $100$ cycles and $7000$ has been added. This can be discarded as we're just looking for the last three digits.Adding up the first $17$ of the cycle of $20$, we can see that the answer is $\boxed{069}$.

 


	\item (\textit{AIME I 2017}) A pyramid has a triangular base with side lengths $20$, $20$, and $24$. The three edges of the pyramid from the three corners of the base to the fourth vertex of the pyramid all have length $25$. The volume of the pyramid is $m\sqrt{n}$, where $m$ and $n$ are positive integers, and $n$ is not divisible by the square of any prime. Find $m+n$.


 
 \textbf{Answer: }803



\textbf{Solution}

Let the triangular base be $\triangle ABC$, with $\overline {AB} = 24$. We find that the altitude to side $\overline {AB}$ is $16$, so the area of $\triangle ABC$ is $(24*16)/2 = 192$.

 Let the fourth vertex of the tetrahedron be $P$, and let the midpoint of $\overline {AB}$ be $M$. Since $P$ is equidistant from $A$, $B$, and $C$, the line through $P$ perpendicular to the plane of $\triangle ABC$ will pass through the circumcenter of $\triangle ABC$, which we will call $O$. Note that $O$ is equidistant from each of $A$, $B$, and $C$. Then,

 
 \[\overline {OM} + \overline {OC} = \overline {CM} = 16\]

 Let $\overline {OM} = d$. Then $OC=OA=\sqrt{d^2+12^2}.$Equation $(1)$:
 \[d + \sqrt {d^2 + 144} = 16\]

 Squaring both sides, we have

 
 \[d^2 + 144 + 2d\sqrt {d^2+144} + d^2 = 256\]

 \[2d^2 + 2d\sqrt {d^2+144} = 112\]

 \[2d(d + \sqrt {d^2+144}) = 112\]

 Substituting with equation $(1)$:

 
 \[2d(16) = 112\]

 \[d = 7/2\]

 We now find that $\sqrt{d^2 + 144} = 25/2$.

 Let the distance $\overline {OP} = h$. Using the Pythagorean Theorem on triangle $AOP$, $BOP$, or $COP$ (all three are congruent by SSS):

 
 \[25^2 = h^2 + (25/2)^2\]

 \[625 = h^2 + 625/4\]

 \[1875/4 = h^2\]

 \[25\sqrt {3} / 2 = h\]

 
Finally, by the formula for volume of a pyramid,

 
 \[V = Bh/3\]

 \[V = (192)(25\sqrt{3}/2)/3\]This simplifies to $V = 800\sqrt {3}$, so $m+n = \boxed {803}$.

 
NOTE : If you don’t know or remember the formula for the volume of a triangular pyramid, you can derive it using calculus as follows :

 Take a small triangular element in the pyramid. We know that it’s area is proportional to the height from the vertex to the base. Hence, we know that $\frac{A_{small element}}{A} = \frac{h^2}{H^2} \implies A_{small element} = \frac{Ah^2}{H^2}$. Now integrate it taking the limits $0$ to $H$

 \textbf{Shortcut} Here is a shortcut for finding the radius $R$ of the circumcenter of $\triangle ABC$.

 As before, we find that the foot of the altitude from $P$ lands on the circumcenter of $\triangle ABC$. Let $BC=a$, $AC=b$, and $AB=c$.Then we write the area of $\triangle ABC$ in two ways:
 \[[ABC]= \frac{1}{2} \cdot 24 \cdot 16 = \frac{abc}{4R}\]

 Plugging in $20$, $20$, and $24$ for $a$, $b$, and $c$ respectively, and solving for $R$, we obtain $R= \frac{25}{2}=OA=OB=OC$.

 Then continue as before to use the Pythagorean Theorem on $\triangle AOP$, find $h$, and find the volume of the pyramid.

 \textbf{Another Shortcut (Extended Law of Sines)} Take the base $\triangle ABC$, where $AB = BC = 20$ and $AC = 24$. Draw an altitude from $B$ to $AC$ that bisects $AC$ at point $D$. Then the altitude has length $\sqrt{20^2 - 12^2} = \sqrt{16^2} = 16$. Next, let $\angle BCA = \theta$. Then from the right triangle $\triangle BDC$, $\sin \theta = 4/5$. From the extended law of sines, the circumradius is $20 \cdot \dfrac{5}{4} \cdot \dfrac{1}{2} = \dfrac{25}{2}$.

 



\textbf{Solution 2 (Coordinates)}

We can place a three dimensional coordinate system on this pyramid. WLOG assume the vertex across from the line that has length $24$ is at the origin, or $(0, 0, 0)$. Then, the two other vertices can be $(-12, -16, 0)$ and $(12, -16, 0)$. Let the fourth vertex have coordinates of $(x, y, z)$. We have the following $3$ equations from the distance formula.

 
 \[x^2+y^2+z^2=625\]

 \[(x+12)^2+(y+16)^2+z^2=625\]

 \[(x-12)^2+(y+16)^2+z^2=625\]

 Adding the last two equations and substituting in the first equation, we get that $y=-\frac{25}{2}$. If you drew a good diagram, it should be obvious that $x=0$. Now, solving for $z$, we get that $z=\frac{25\sqrt{3}}{2}$. So, the height of the pyramid is $\frac{25\sqrt{3}}{2}$. The base is equal to the area of the triangle, which is $\frac{1}{2} \cdot 24 \cdot 16 = 192$. The volume is $\frac{1}{3} \cdot 192 \cdot \frac{25\sqrt{3}}{2} = 800\sqrt{3}$. Thus, the answer is $800+3 = \boxed{803}$.

 \textbf{-RootThreeOverTwo}

 



\textbf{Solution 3 (Heron's Formula)}

Label the four vertices of the tetrahedron and the midpoint of $\overline {AB}$, and notice that the area of the base of the tetrahedron, $\triangle ABC$, equals $192$, according to Solution 1. 

 Notice that the altitude of $\triangle CPM$ from $\overline {CM}$ to point $P$ is the height of the tetrahedron. Side $\overline {PM}$ is can be found using the Pythagorean Theorem on $\triangle APM$, giving us $\overline {PM}=\sqrt{481}.$ 

 Using Heron's Formula, the area of $\triangle CPM$ can be written as 
 \[\sqrt{\frac{41+\sqrt{481}}{2}(\frac{41+\sqrt{481}}{2}-16)(\frac{41+\sqrt{481}}{2}-25)(\frac{41+\sqrt{481}}{2}-\sqrt{481})}\]
 \[=\frac{\sqrt{(41+\sqrt{481})(9+\sqrt{481})(-9+\sqrt{481})(41-\sqrt{481})}}{4}\]

 Notice that both $(41+\sqrt{481})(41-\sqrt{481})$ and $(9+\sqrt{481})(-9+\sqrt{481})$ can be rewritten as differences of squares; thus, the expression can be written as 
 \[\frac{\sqrt{(41^2-481)(481-9^2)}}{4}=\frac{\sqrt{480000}}{4}=100\sqrt{3}.\]

 From this, we can determine the height of both $\triangle CPM$ and tetrahedron $ABCP$ to be $\frac{100\sqrt{3}}{8}$; therefore, the volume of the tetrahedron equals $\frac{100\sqrt{3}}{8} \cdot 192=800\sqrt{3}$; thus, $m+n=800+3=\boxed{803}.$

 \textbf{-dzhou100}

 


 



\textbf{Solution 4 (Symmetry)}


                    \begin{figure}[H]
                        \centering
                        \includegraphics[width=0.5\textwidth]{Solution/AIME_I_2017_4_Solution_Figure_1.png}
                    \end{figure}
                    Notation is shown on diagram. 
 \[AM = MB = c = 12, AC = BC = b = 20,\]
 \[DA = DB = DC = a = 25.\]
 \[CM = x + y = \sqrt{b^2-c^2} = 16,\] 
 \[x^2 - y^2 = CD^2 – DM^2 = CD^2 – (BD^2 – BM^2) = c^2 = 144,\]
 \[x – y = \frac{x^2 – y^2}{x+y} = \frac {c^2} {16} = 9,\]
 \[x = \frac {16 + 9}{2} = \frac {a}{2},\]
 \[h = \sqrt{a^2 -\frac{ a^2}{4}} = a \frac {\sqrt{3}}{2},\]
 \[V = \frac{h\cdot CM \cdot c}{3}= \frac{16\cdot 25 \sqrt{3} \cdot 12}{3} = 800 \sqrt{3} \implies \boxed {803}.\]\textbf{vladimir.shelomovskii@gmail.com, vvsss}

 



\textbf{Solution 5 (Similar to others)}

\textbf{Fact 1:}

 $D$ is the circumsphere of $ABC$.

 \textbf{Lemma:}

 The altitude of the tetrahedron from apex $D$ intersects plane $ABC$ at the circumcentre.

 \textbf{Proof of Lemma:}

 The horizontal plane that passes through any part of the sphere always maps a perfect circle, so the horizontal plane that passes through triangle $ABC$ passes through a circle. The diameter of the sphere perpendicular to the plane $ABC$ has the property that it will always be located at the centre of any given plane. Therefore, say this line intersects $ABC$ at $O$, then $O$ is the centre of the circle that passes through $ABC$, and therefore the circumcentre.

 \textbf{2017 AIME I Q4:}

 Through simple formulae the circumradius is $\frac{25}{2}$. Therefore, the triangle containing $AOD$ is 30-60-90 and $DO = \frac{25\sqrt{3}}{2}$. From there, the area of the base triangle $ABC$ is 192, so we find $32 \cdot 25\sqrt{3} = 800\sqrt{3} \implies \boxed{803}$.

 


	\item (\textit{AIME I 2017}) A rational number written in base eight is $\underline{a} \underline{b} . \underline{c} \underline{d}$, where all digits are nonzero. The same number in base twelve is $\underline{b} \underline{b} . \underline{b} \underline{a}$. Find the base-ten number $\underline{a} \underline{b} \underline{c}$.


 
 \textbf{Answer: }321



\textbf{Solution 1}

First, note that the first two digits will always be a positive number. We will start with base twelve because of its repetition. List all the positive numbers in base eight that have equal tens and ones digits in base 12.

 $11_{12}=15_8$
$22_{12}=32_8$
$33_{12}=47_8$
$44_{12}=64_8$
$55_{12}=101_8$ 

 We stop because we only can have two-digit numbers in base 8 and 101 is not a 2 digit number.Compare the ones places to check if they are equal. We find that they are equal if $b=2$ or $b=4$.Evaluating the places to the right side of the decimal point gives us $22.23_{12}$ or $44.46_{12}$.When the numbers are converted into base 8, we get $32.14_8$ and $64.30_8$. Since $d\neq0$, the first value is correct. Compiling the necessary digits leaves us a final answer of $\boxed{321}$




\textbf{Solution 2}

The parts before and after the decimal points must be equal. Therefore $8a + b = 12b + b$ and $c/8 + d/64 = b/12 + a/144$.  Simplifying the first equation gives $a = (3/2)b$.  Plugging this into the second equation gives $3b/32 = c/8 + d/64$. Multiplying both sides by 64 gives $6b = 8c + d$.  $a$ and $b$ are both digits between 1 and 7 (they must be a single non-zero digit in base eight) so using $a = 3/2b$, $(a,b) = (3,2)$ or $(6,4)$.  Testing these gives that $(6,4)$ doesn't work, and $(3,2)$ gives $a = 3, b = 2, c = 1$, and $d = 4$.  Therefore $abc = \boxed{321}$




\textbf{Solution 3}

Converting to base $10$ we get 

 $4604a+72c+9d=6960b$

 Since $72c$ and $9d$ are much smaller than the other two terms, dividing by $100$ and approximating we get 

 $46a=70b$

 Writing out the first few values of $a$ and $b$, the first possible tuple is 

 $a=3, b=2, c=1, d=4$ 

 and the second possible tuple is

 $a=6, b=4, c=3, d=0$

 Note that $d$ can not be $0$, therefore the answer is $\boxed{321}$

 By maxamc

 


 


 Note: Another way to get the tuples is by setting the different numbers equal to each other.

 Doing that we get:

 $8a + b + \frac{c}{8} + \frac{d}{64} = 12b + b + \frac{b}{12} + \frac{a}{144}$

 Moving the integer terms to one side we get:

 $8a - 12b = \frac{b}{12} + \frac{a}{144} - (\frac{c}{8} + \frac{d}{64})$

 Factoring the LHS:

 $4(2a - 3b) = \frac{b}{12} + \frac{a}{144} - (\frac{c}{8} + \frac{d}{64})$

 Since $a$ and $b$ are integers the RHS must also be a Integer and therefore a multiple of 4. However, due to the bounding of $a, b, c, d < 10$ the RHS must be 0, as it cannot reach 4 or -4.

 Thus, $2a - 3b = 0$ and $2a = 3b$

 From here you can bash and find there are only 2 possible tuples.

 



\textbf{Solution 4 (modular congruency/nt bash)}

In the problem, we are given that 
 \[8a+b+\dfrac c8+\dfrac d{64}=12b+b+\dfrac b{12}+\dfrac a{144}.\] We multiply by the LCM of the denominators, which is $576$ to get 
 \[4608a+576b+72c+9d=6912b+576b+48b+4a.\] We then group like terms and factor to get 
 \[4604a+72c+9d=48b(144+1)=145\cdot48b=24\cdot290b.\] 

 Observe that the coefficients of $a$, $c$, and $b$ are all divisible by $4$. Therefore, we know that $d$ must also be divisible by $4$ to compensate. (To observe this, one could rearrange the terms to see $9d=24\cdot290b-4604a-72c=4(6\cdot290b-1151a-18c)$. At this point, it is obvious that $9d$ must be divisible by $4$, so $d$ must be divisible by $4$.) Thus, we let $d=4y$ to see $4604a+72c+9\cdot4y=24\cdot290b$. We divide by $4$ in the equation to get 
 \[1151a+18c+9y=6\cdot290b.\] 

 Observe that the coefficients of $c$, $y$, and $b$ are all divisible by $3$. By similar reasoning, $a$ must be divisible by $3$, so we let $a=3x$. Substituting and dividing by $3$, we get 
 \[1151x+6c+3y=2\cdot290b.\] We observe that this can no longer be reduced by similar means. 

 We know that $0<a,b,c<8\implies0<3x,b,4y<8$. We examine $0<4y<8$; this becomes $0<y<2$. It is apparent that $y=1\implies d=4$. However, the problem does not even ask for $d$, so it may appear that this find is meaningless. However, we substitute our value of $y$ in to get 
 \[1151x+6c+3=580b.\] 

 We know that $x$ is either $1$ or $2$, since $0<3x<8$. We take our equation modulo $6$ to find 
 \[5x+3\equiv4b\pmod{6}\implies3\equiv4b+x\pmod{6}.\] If $x=2$, then $1\equiv4b\pmod{6}$ - this is equivalent to saying that $1+6n=4b\implies 1=4b-6n$. However, $4b-6n$ is always even, and $1$ is odd. Thus, this case is not possible, so $x=1\implies a=3$. We know that $x=1$; thus, $3\equiv4b+1\pmod{6}\implies2\equiv4b\pmod{6}\implies1\equiv2b\pmod{3}$. Obviously, $b=2$ and $b=5$ work; however, if $b=5$, then the RHS of our original equation ($2\cdot290b$) is much too large to be equal to the LHS (the maximum possible value of the LHS is $1151\cdot1+6\cdot7+3\cdot1$, which is less than $1200$ while the RHS would become $2\cdot290\cdot5=2900$), so we have $b=2$. 

 We recall our equation of $1151x+6c+3=580b$. Plugging in what we know, we have $1151\cdot1+6c+3=580\cdot2\implies6c=6\implies c=1$. 

 Therefore, $\overline{abc}=\boxed{321}$. 

 by Technodoggo

 


	\item (\textit{AIME I 2017}) A circle circumscribes an isosceles triangle whose two congruent angles have degree measure $x$. Two points are chosen independently and uniformly at random on the circle, and a chord is drawn between them. The probability that the chord intersects the triangle is $\frac{14}{25}$. Find the difference between the largest and smallest possible values of $x$.


 
 \textbf{Answer: }048



\textbf{Solution}

The probability that the chord doesn't intersect the triangle is $\frac{11}{25}$. The only way this can happen is if the two points are chosen on the same arc between two of the triangle vertices. The probability that a point is chosen on one of the arcs opposite one of the base angles is $\frac{2x}{360}=\frac{x}{180}$ (this comes from the Central Angle Theorem, which states that the central angle from two points on a circle is always twice the inscribed angle from those two points), and the probability that a point is chosen on the arc between the two base angles is $\frac{180-2x}{180}$. Therefore, we can write 
 \[2\left(\frac{x}{180}\right)^2+\left(\frac{180-2x}{180}\right)^2=\frac{11}{25}\]This simplifies to 
 \[x^2-120x+3024=0\]

 (Note that the simplification is quite tedious)

 Which factors as 
 \[(x-84)(x-36)=0\]So $x=84, 36$. The difference between these is $\boxed{048}$.

 
Note:

 We actually do not need to spend time factoring $x^2 - 120x + 3024$. Since the problem asks for $|x_1 - x_2|$, where $x_1$ and $x_2$ are the roots of the quadratic, we can utilize Vieta's by noting that $(x_1 - x_2) ^ 2 = (x_1 + x_2) ^ 2 - 4x_1x_2$. Vieta's gives us $x_1 + x_2 = 120,$ and $x_1x_2 = 3024.$ Plugging this into the above equation and simplifying gives us $(x_1 - x_2) ^ 2 = 2304,$ or $|x_1 - x_2| = 48$. 

 Our answer is then $\boxed{048}$.

 Another note: 

 Letting $y = x/180$, the first equation turns into $2y^2 + (1-2y)^2 = 11/25$, heavily simplifying the otherwise messy simplification, and you can use the aforementioned note to find $|y_1 - y_2| = 4/15$. Using $x = 180y$, $|x_1 - x_2| =  180|y_1 - y_2|$, so the answer is $180(4/15) = \boxed{048}$




\textbf{Solution 2 (Not Complementary Counting method)}

Because we know that we have an isosceles triangle with angles of $x$ (and we know that x is an inscribed angle), that means that the arc that is intercepted by this angle is $2x$. We form this same conclusion for the other angle $x$, and $180-2x$. Therefore we get $3$ arcs, namely, $2x$, $2x$, and $360-4x$. To have the chords intersect the triangle, we need the two points selected (to make a chord) to be on completely different arcs. An important idea to understand is that order matters in this case, so we have the equation $2$ * $\frac{2x}{360}$ * $\frac{2x}{360}$ + $2$ * $2$ * $\frac{2x}{360}$ * $\frac{360-4x}{360}$ = $\frac{14}{25}$ which using trivial algebra gives you $x^2-120x+3024$ and factoring gives you $(x-84)(x-36)$ and so your answer is $\boxed{048}$. ~jske25

 



\textbf{Solution 3 ( 3 System Algebra )}

After constructing the circumscribed circle, realize that the only time when the chord does not intersect the circle is when our $2$ points fall on only one arc formed by the sides of the triangle. Thus, lets call our isosceles triangle  $ABC$, where $AB=BC$. Thus, the arcs formed by $BC$ and $AB$ can be called $a$, and the arc formed by $AC$ is called $b$. So, we can create the following system

 
 \[2a+b=1\]
 \[2a^2+4ab=\frac{14}{25}\]
 \[2a^2+b^2=\frac{11}{25}\]

 Notice we are denoting $a$ and $b$ as our probabilities, which we will be converting to degrees later. The 2 remaining systems can be calculated by using our rule about intersecting arcs and chords. So, after some hairy algebra we get: 

 
 \[a=\frac{1}{5}\] if 
 \[b=\frac{3}{5}\]
 \[a=\frac{7}{15}\] if 
 \[b=\frac{1}{15}\]

 From here we find the absolute difference by doing $\frac{7}{15}-\frac{1}{5} = \frac{4}{15}$. Converting to degrees, since the angles of a triangle add up to $180^o$, we find that $\frac{4}{15} \cdot 180 =\boxed{048}$, which is our answer.

 -Geometry285

 


	\item (\textit{AIME I 2017}) For nonnegative integers $a$ and $b$ with  $a + b \leq 6$, let $T(a, b) = \binom{6}{a} \binom{6}{b} \binom{6}{a + b}$. Let $S$ denote the sum of all $T(a, b)$, where  $a$ and $b$ are nonnegative integers with $a + b \leq 6$. Find the remainder when $S$ is divided by $1000$.


 
 \textbf{Answer: }564



\textbf{Solution 1 (Committee Forming)}

Let $c=6-(a+b)$, and note that $\binom{6}{a + b}=\binom{6}{c}$. The problem thus asks for the sum $\binom{6}{a} \binom{6}{b} \binom{6}{c}$ over all $a,b,c$ such that $a+b+c=6$. Consider an array of 18 dots, with 3 columns of 6 dots each. The desired expression counts the total number of ways to select 6 dots by considering each column separately, which is equal to $\binom{18}{6}=18564$. Therefore, the answer is $\boxed{564}$.

 -rocketscience

 



\textbf{Solution 1 but different (Committee Forming)}

Alternatively, one can note that we can consider groups where $a+b$ is constant, say $c$. Fix any value of $c$. Then the sum of all of the values of $T(a,b)$ such that $a+b=c$ is $\binom{6}{a+b} \sum_{a+b=c} \binom{6}{a}\binom{6}{b}$ which by Vandermonde's is $\binom{6}{a+b}\binom{12}{a+b}$. Remember, that expression is the resulting sum for a fixed $a+b$. So, for $a+b\le 6$, we want $\sum_{c=0}^{6} \binom{6}{c}\binom{12}{c}$. This is (by Vandermonde's or committee forming) $\binom{18}{6} = 18564 \implies \boxed{564}$ ~ firebolt360

 \textbf{Note} Now just a quick explanation for people who don't fully understand Vandermonde's. Take the first part, $\sum_{a+b=c} \binom{6}{a}\binom{6}{b}$. Consider $2$ different groups, $A$ and $B$ both of size $6$ people. We wish to chose $a$ peoples from $A$ and $b=c-a$ people from $B$. In total, we chose $c-a+a=c$ people. We can then draw a bijection towards choosing $c$ people from $A\cup B$, which has size $12$. So, it is $\binom{12}{c}=\binom{12}{a+b}$. Similarly, for $\sum_{c=6} \binom{6}{c}\binom{12}{c}$, we see that $\binom{6}{c}=\binom{6}{6-c}$. Now the total is $18$, and the sum is $6$. So, we get $\binom{18}{6}$. See \href{https://artofproblemsolving.com/wiki/index.php?title=Committee\_forming}{committee forming} for more information ~ firebolt360

 



\textbf{Solution 2 (Committee Forming but slightly more bashy)}

Treating $a+b$ as $n$, this problem asks for 
 \[\sum_{n=0}^{6} \left[\binom{6}{n}\sum_{m=0}^{n}\left[\binom{6}{m}\binom{6}{n-m}\right]\right].\] But 
 \[\sum_{m=0}^{n} \left[\binom{6}{m} \binom{6}{n-m}\right]\]can be computed through the following combinatorial argument.  Choosing $n$ elements from a set of size $12$ is the same as splitting the set into two sets of size $6$ and choosing $m$ elements from one, $n-m$ from the other where $0\leq m\leq n$.  The number of ways to perform such a procedure is simply $\binom{12}{n}$. Therefore, the requested sum is 
 \[\sum_{n=0}^{6} \left[\binom{6}{n} \binom{12}{n}\right] = 18564.\] As such, our answer is $\boxed{564}$.

 Note from epiconan:To avoid the bash, you can actually use Vandermondes $\emph{again}$.
 \[\sum_{n=0}^{6} \left[\binom{6}{n} \binom{12}{n}\right] = \sum_{n=0}^{6} \left[\binom{6}{6-n} \binom{12}{n}\right] = \binom{6+12}{6} = \binom{18}{6}.\]- Awsomness2000

 



\textbf{Solution 3 (Major Major Bash)}

Case 1: $a<b$.

 Subcase 1: $a=0$
 \[\binom{6}{0}\binom{6}{1}\binom{6}{1}=36\]
 \[\binom{6}{0}\binom{6}{2}\binom{6}{2}=225\]
 \[\binom{6}{0}\binom{6}{3}\binom{6}{3}=400\]
 \[\binom{6}{0}\binom{6}{4}\binom{6}{4}=225\]
 \[\binom{6}{0}\binom{6}{5}\binom{6}{5}=36\]
 \[\binom{6}{0}\binom{6}{6}\binom{6}{6}=1\]
 \[36+225+400+225+36+1=923\]Subcase 2: $a=1$
 \[\binom{6}{1}\binom{6}{2}\binom{6}{3}=1800 \equiv 800 \pmod {1000}\]
 \[\binom{6}{1}\binom{6}{3}\binom{6}{4}=1800 \equiv 800 \pmod {1000}\]
 \[\binom{6}{1}\binom{6}{4}\binom{6}{5}=540\]
 \[\binom{6}{1}\binom{6}{5}\binom{6}{6}=36\]
 \[800+800+540+36=2176 \equiv 176 \pmod {1000}\]Subcase 3: $a=2$
 \[\binom{6}{2}\binom{6}{3}\binom{6}{5}=1800\equiv800\pmod{1000}\]
 \[\binom{6}{2}\binom{6}{4}\binom{6}{6}=225\]
 \[800+225=1025\equiv25\pmod{1000}\]

 \[923+176+25=1124\equiv124\pmod{1000}\]

 Case 2: $b<a$

 By just switching $a$ and $b$ in all of the above cases, we will get all of the cases such that $b>a$ is true. Therefore, this case is also $124\pmod{1000}$

 Case 3: $a=b$
 \[\binom{6}{0}\binom{6}{0}\binom{6}{0}=1\]
 \[\binom{6}{1}\binom{6}{1}\binom{6}{2}=540\]
 \[\binom{6}{2}\binom{6}{2}\binom{6}{4}=3375\equiv375\pmod{1000}\]
 \[\binom{6}{3}\binom{6}{3}\binom{6}{6}=400\]
 \[1+540+375+400=1316\equiv316\pmod{1000}\]

 \[316+124+124=\boxed{\boxed{564}}\]




\textbf{Solution 4}

We begin as in solution 1 to rewrite the sum as $\binom{6}{a} \binom{6}{b} \binom{6}{c}$ over all $a,b,c$ such that $a+b+c=6$.Consider the polynomial $P(x)=\left(\binom{6}{0} + \binom{6}{1}x + \binom{6}{2}x^2 + \binom{6}{3}x^3 + \cdot \cdot \cdot + \binom{6}{6}x^6\right)^3$.We can see the sum we wish to compute is just the coefficient of the $x^6$ term. However $P(x)=((1+x)^6)^3=(1+x)^{18}$. Therefore, the coefficient of the $x^6$ term is just $\binom{18}{6} = 18564$ so the answer is $\boxed{564}$.

 - mathymath

 



\textbf{Solution 5 (Committee Forming but different)}

Let $c=6-(a+b)$. Then $\binom{6}{a+b}=\binom{6}{c}$, and $a+b+c=6$. The problem thus asks for 
 \[\sum_{a+b+c=6}\binom{6}{a}\binom{6}{b}\binom{6}{c} \pmod {1000}.\] Suppose we have $6$ red balls, $6$ green balls, and $6$ blue balls lined up in a row, and we want to choose $6$ balls from this set of $18$ balls by considering each color separately. Over all possible selections of $6$ balls from this set, there are always a nonnegative number of balls in each color group. The answer is $\binom{18}{6} \pmod {1000}=18\boxed{564}$.

 



\textbf{Solution 5 but different (Committee Forming)}

Since $\binom{6}{n}=\binom{6}{6-n}$, we can rewrite $T(a,b)$ as $\binom{6}{a}\binom{6}{b}\binom{6}{6-(a+b)}$. Consider the number of ways to choose a committee of 6 people from a group of 6 democrats, 6 republicans, and 6 independents. We can first pick $a$ democrats, then pick $b$ republicans, provided that $a+b \leq 6$. Then we can pick the remaining $6-(a+b)$ people from the independents. But this is just $T(a,b)$, so the sum of all $T(a,b)$ is equal to the number of ways to choose this committee.On the other hand, we can simply pick any 6 people from the $6+6+6=18$ total politicians in the group. Clearly, there are $\binom{18}{6}$ ways to do this. So the desired quantity is equal to $\binom{18}{6}$. We can then compute (routinely) the last 3 digits of $\binom{18}{6}$ as $\boxed{564}$.

 



\textbf{Solution 6(NICE Journal)}

Note that $\binom{6}{a+b} = \binom{6}{6-a-b}$. So we have $\binom{6}{a}\binom{6}{b}\binom{6}{6-a-b}$. If we think about this this is essentially choosing a group of $a$ people from $6$ people, a group of $b$ people from $6$ people, and a group of $6 - a -b$ from another group of $6$ people. This is nothing but choosing $6$ people from a group of $18$ people. This is nothing but $\binom{18}{6} = 18564 \Rightarrow 564$.~coolmath\_2018

 



\textbf{Solution 7 (The Bash of 2026)}

Obviously, one can do this on an AIME, this just took me 20 mins, mainly because im good at computations. If you suck at calculating fast, then just try learning committee forming, it is much easier than this.

 Case 1: $a+b = 0$
$a = 0, b = 0$ so $\binom{6}{0} \binom{6}{0} \binom{6}{0} = 1$.

 Running total 1.

 Case 2: $a+b = 1$
$a = 0, b = 1$ or vice versa so $2 \cdot \binom{6}{0} \binom{6}{1} \binom{6}{1} = 72$

 Running total 73.

 Case 3: $a+b = 2$
$a = 0, b = 2$ or vice versa, so $2 \cdot \binom{6}{0} \binom{6}{2} \binom{6}{2} = 450$
$a = 1, b=1$ so $\binom{6}{1} \binom{6}{1} \binom{6}{2} = 36 \cdot 15 = 540$.

 Running total 1063, take 63.

 Case 4: $a+b = 3$
$a = 0, b = 3$ or vice versa so $2 \cdot \binom{6}{3} \binom{6}{3} = 800$
$a = 1, b = 2$ or vice versa so $2 \cdot \binom{6}{1} \binom{6}{2} \binom{6}{3} = 2 \cdot 6 \cdot 15 \cdot 20 = 3600$.

 Running total 4463, take 463.

 Case 5: $a+b = 4$
$a = 0, b = 4$ or vice versa, so $2 \cdot \binom{6}{4} \cdot \binom{6}{4} = 450$.

 $a = 1, b = 3$ or vice versa so $2 \cdot \binom{6}{1} \binom{6}{3} \binom{6}{4} = 3600$
$a = 2, b = 2$ so $15 \cdot 15 \cdot 15 = 3375$.

 Running total 7888. Take 888.

 Case 6: $a+b = 5$
$a= 0, b  = 5$ or vice versa, so $2 \cdot 6 \cdot 1 \cdot 6 = 72$.

 $a=1,b=4$ or vice versa so $2 \cdot 6 \cdot 15 \cdot 6  = 1080$.

 $a=2, b=3$ or vice versa so $2 \cdot 15 \cdot 20 \cdot 6 = 3600$.

 Running total 5640, take 640.

 Case 7: $a+b = 6$
$a = 0, b = 6$ or vice versa so $2 \cdot 1 \cdot 1 \cdot 1 = 2$
$a=1, b= 5$ or vice versa so $2 \cdot 6 \cdot 6 \cdot 1 = 72$
$a= 2, b=4$ or vice versa so $2 \cdot 15 \cdot 15 \cdot 1 = 450$
$a = 3, b = 3$ so $20 \cdot 20 = 400$.

 Running total 1564, take $\boxed{564}$.

 



\textbf{Solution 8 (Bijectivity)}

Consider set $A = \{1,2,3,4,5,6\}$, set $B = \{7,8,9,10,11,12 \}$. Define $S(a,b) = \binom{6}{a} \cdot \binom{6}{b}$, so $T(a,b) = S(a,b) \cdot \binom{6}{a+b}$. Define $C = \{13,14,15,16,17,18 \}$. Literally, we choose $a$ elements from $A$, $b$ from $B$, and then $a+b$ elements from $C$, conjoin them to this one big set $A_a \cup B_b \cup C_{a+b}$ of $2a+2b$ elements. Consider $A \cup B \cup C$. Then, out of 18 such elements we literally choose at maximum 12 of them, at minimum 0, so there are $\binom{18}{0} + \binom{18}{2} + \binom{18}{4} + \dots + \binom{18}{12}$ such sets. Now the problem does NOT ask for the number of sets. It asks for $\sum T(a,b)$, which can be obtained by any subset of $\binom{18}{12}$ elements, as such a subset can yield all plausible cases for the remainder of the sets. So, the requested value is $\binom{18}{12} \pmod{1000} \equiv \boxed{564}$.

 
 \[\textbf{OR}\]

 Continue upon realization that $2a+2b$ elements is required in total for the one giant set. Consider $A \cup B \cup C$. Suppose the empty set is not allowed. Then, out of 18 such elements we literally choose $2(a+b)$, which yields 6 different, unique values. Through principle of "Committee Forming", one can correspond all such cases to the sum through simple basis. Therefore, we have $\binom{18}{6} \equiv \boxed{564} \pmod{1000}$ as the sum.

 \textbf{Remark} This problem is an example of the generalization of Vandermonde's theorem, which states that for nonnegative $k_1, k_2, \ldots k_p$ and $n_1, n_2, \ldots n_p$, we have
 \[\sum _{k_1+\cdots +k_p=m}\binom{n_1}{k_1} \binom{n_2}{k_2} \cdots \binom{n_p}{k_p} = \binom{n_1+ \cdots +n_p}{m}.\]~eibc

 


	\item (\textit{AIME I 2017}) Two real numbers $a$ and $b$ are chosen independently and uniformly at random from the interval $(0, 75)$. Let $O$ and $P$ be two points on the plane with $OP = 200$. Let $Q$ and $R$ be on the same side of line $OP$ such that the degree measures of $\angle POQ$ and $\angle POR$ are $a$ and $b$ respectively, and $\angle OQP$ and $\angle ORP$ are both right angles. The probability that $QR \leq 100$ is equal to $\frac{m}{n}$, where $m$ and $n$ are relatively prime positive integers. Find $m + n$.


 
 \textbf{Answer: }041



\textbf{Solution 1}

Noting that $\angle OQP$ and $\angle ORP$ are right angles, we realize that we can draw a semicircle with diameter $\overline{OP}$ and points $Q$ and $R$ on the semicircle. Since the radius of the semicircle is $100$, if $\overline{QR} \leq 100$, then $\overarc{QR}$ must be less than or equal to $60^{\circ}$.

 This simplifies the problem greatly. Since the degree measure of an angle on a circle is simply half the degree measure of its subtended arc, the problem is simply asking:

 Given $a, b$ such that $0<a, b<75$, what is the probability that $|a-b| \leq 30$? Through simple geometric probability, we get that $P = \frac{16}{25}$.

 The answer is $16+25=\boxed{041}$




\textbf{Solution 2 (Trig Bash)}

Put $\triangle POQ$ and $\triangle POR$ with $O$ on the origin and the triangles on the $1^{st}$ quadrant.The coordinates of $Q$ and $P$ is $(200 \cos^{2}a,200 \cos a\sin a )$, $(200\cos^{2}b,200\cos(b)\sin b)$. So $PQ^{2}$ = $(200 \cos^{2} a - 200 \cos^{2} b)^{2} +(200 \cos a \sin a - 200 \cos b \sin b)^{2}$, which we want to be less then $100^{2}$.So $(200 \cos^{2} a - 200 \cos^{2} b)^{2} +(200 \cos a \sin a - 200 \cos b \sin b)^{2} \le 100^{2}$
 \[(\cos^{2} a - \cos^{2} b)^{2} +(\cos a \sin a - \cos b \sin b)^{2} \le \frac{1}{4}\]
 \[\cos^{4} a + \cos^{4} b - 2\cos^{2} a \cos^{2} b +\cos^{2}a \sin^{2} a + \cos^{2} b \sin^{2} b - 2 \cos a \sin a \cos b \sin b \le \frac{1}{4}\]
 \[\cos^{2} a(\cos^{2} a + \sin^{2} a)+\cos^{2} b(\cos^{2} b+\sin^{2} b) - 2\cos^{2} a \cos^{2} b- 2 \cos a \sin a \cos b \sin b \le \frac{1}{4}\]
 \[\cos^{2} a(1-\cos^{2} b)+\cos^{2} b(1-\cos^{2} a) - 2 \cos a \sin a \cos b \sin b \le \frac{1}{4}\]
 \[(\cos a\sin b)^{2} +(\cos b\sin a)^{2} - 2 (\cos a \sin b)(\cos b \sin a)\le \frac{1}{4}\]
 \[(\cos a\sin b-\cos b\sin a)^{2}\le \frac{1}{4}\]
 \[\sin^{2} (b-a) \le \frac{1}{4}\]So we want $-\frac{1}{2} \le \sin (b-a) \le \frac{1}{2}$, which is equivalent to $-30 \le b-a \le 30$ or $150 \le b-a \le 210$. The second inequality is impossible so we only consider what the first inequality does to our $75$ by $75$ box in the $ab$ plane. This cuts off two isosceles right triangles from opposite corners with side lengths $45$ from the $75$ by $75$ box. Hence the probability is $1-\frac{45^2}{75^2} = 1- \frac{9}{25}=\frac{16}{25}$ and the answer is $16+25 = \boxed{41}$

 Solution by Leesisi

 



\textbf{Solution 3 (Quicker Trig)}


                    \begin{figure}[H]
                        \centering
                        \includegraphics[width=0.5\textwidth]{Solution/AIME_I_2017_8_Solution_Figure_1.png}
                    \end{figure}
                    Let $QR=x.$ Since we are given many angles in the problem, we can compute the lengths of some of the lines in terms of trigonometric functions: $OQ = 200 \cos a, PQ = 200 \sin a, PR = 200 \sin b, OR = 200 \cos b.$ Now observe that quadrilateral $OQRP$ is a \href{https://artofproblemsolving.com/wiki/index.php?title=Cyclic\_quadrilateral}{cyclic quadrilateral}. Thus, we are able to apply \href{https://artofproblemsolving.com/wiki/index.php?title=Ptolemy\%27s\_Theorem}{Ptolemy's Theorem} to it:
 \[200 x + (200 \cos a) (200 \sin b) = (200 \sin a) (200 \cos b),\]
 \[x + 200 (\cos a \sin b) = 200 (\sin a \cos b),\]
 \[x = 200(\sin a \cos b - \sin b \cos a),\]
 \[x = 200 \sin(a-b).\]We want $|x| \le 100$ (the absolute value comes from the fact that $a$ is not necessarily greater than $b,$ so we cannot assume that $Q$ is to the right of $R$ as in the diagram), so we substitute:
 \[|200 \sin(a-b)| \le 100,\]
 \[|\sin(a-b)| \le \frac{1}{2},\]
 \[|a-b| \le 30 ^\circ,\]
 \[-30 \le a-b \le 30.\]By simple geometric probability (see Solution 2 for complete explanation), $\frac{m}{n} = 1 - \frac{2025}{5625} = 1 - \frac{9}{25} = \frac{16}{25},$ so $m+n = \boxed{041}.$




\textbf{Solution 4}

Scale the circle down from radius 100 (diameter 200) to radius 6 (diameter 12). Then we want the probability that $PQ \le 6$. Now note that all possible $P$ and $Q$ lie on a $5\pi$ interval on the circumference of the circle. But for $PQ<6$, $P$ and $Q$ must be less than $2\pi$ apart on the circumference of the circle. Simple geometric probability gives us $\frac{16}{25}$, so the answer is $\fbox{41}$. (Professor-Mom)

 



\textbf{Solution 5}

Impose a coordinate system as follows:

 Let the midpoint of $\overline{OP}$ be the origin, and let $\overline{OP}$ be the x-axis. We construct a circle with center at the origin with radius 100. Since $\angle OQP$ and $\angle ORP$ are both right angles, points $Q$ and $R$ are on our circle. Place $Q$ and $R$ in the first quadrant of the Cartesian Plane. Suppose we construct $Q'$ and $R'$ such that they are clockwise rotations of $Q$ and $R$, respectively by an angle of $2b$ degrees. Thus, we see that $\overline{QR}=100\sqrt{2}\sqrt{\cos(2|a-b|)}$. We want this quantity to be less than $100$. This happens when $\cos(2|a-b|) \ge 1/2,$ or when $|a-b|\le 30^{\circ}$. The probability that the last inequality is satisfied is $16/25$. Therefore, the probability that $QR$ is less than $100$ is $16/25$. Hence, $m+n=\boxed{41}$

 
slightly edited

 


 



\textbf{Solution 6}

WLOG, let $b\ge a$. It does not actually matter, but it is necessary for this particular setup. It should be apparent that $\Delta RAQ\sim\Delta OAP$. We write the equation 

 
 \[\dfrac{RA}{AO}=\dfrac{RQ}{OP}.\] 

 If we examine right triangle $\Delta ROA$, we can see that $\sin(b-a)=\dfrac{RA}{AO}$. Also, we are given $OP=200$, so now we have 

 
 \[\sin(b-a)=\dfrac{QR}{200}.\]

 We want $QR$ to be less than or equal to $100$; this is equivalent to $\dfrac{QR}{200}\le\dfrac12.$ We solve from there: 

 \begin{align*}\dfrac{QR}{200}&\le\dfrac12 \\\sin(b-a)&\le\dfrac12 \\\arcsin(\sin(b-a))&\le\arcsin\left(\dfrac12\right) \\b-a&\le30^\circ. \\\end{align*}


 (Notice that if $a>b$, then this would become $a-b\le30^\circ.$ As in Solution 1, we can write $|a-b|\le30$.) One can now proceed as in Solution 1, but let us tackle the geometric probability for completeness. 

 We now have transformed this problem into another problem asking for the probability of two uniformly, randomly, and independently chosen real numbers between $0$ and $75$ being no more than $30$ from each other. 

 If the first number (let this be $x$) is between $30$ and $45$, then the other number can be from $x-30$ to $x+30$ - a range of $60$. Thus, the probability that this contributes is $\dfrac{45-30}{75}\cdot\dfrac{60}{75}=\dfrac4{25}$. 

 If $x$ is between $0$ and $30$ or $45$ to $75$ (these two cases are equivalent), the chance is the same as that of the average value since the ranges are uniform. For $x=15$ (the average), the second number can be from $0$ to $15+30=45$ - a range of $45$. The total range is $(30-0)+(75-45)=30+30=60$. Thus, this case contributes $\dfrac{45}{75}\cdot\dfrac{60}{75}=\dfrac{12}{25}$. 

 Adding the two, we get $\dfrac{16}{25}$ for an answer of $16+25=\boxed{041}$. 

 



\textbf{Solution 7}

We notice that the arc $QR=2a-2b$. We let $M$ be the midpoint of $OP$ and the center of the semicircle shown in the diagram in Solution 3. From arc $QR$, we find $\angle QMR=2a-2b$. Because the radius of the semicircle is 100, we can use law of cosines to find the length of $QR$ giving us, 
 \[QR^2=100^2+100^2-2(100)(100)\cos(2a-2b).\]Because $QR<100$, we know $QR^2<100^2$ giving us, 
 \[100^2\ge100^2+100^2-2(100^2)\cos(2a-2b),\]
 \[1\ge2-2\cos(2a-2b),\]
 \[\cos(2a-2b)\ge\frac{1}{2}.\] 

 Because $\cos^{-1}\left(\frac{1}{2}\right)=\pm60^{\circ}$, we know that $2a-2b\le60$ or $2a-2b\ge300$. However, the maximum value of $a$ is $65$ meaning $2a-2b\ge300$ is impossible. Thus, we find, 
 \[a-b<30.\] We can use geometric probability to find the probability of this occurring which ends up giving us the following equation, 
 \[\frac{75^2-45^2}{75^2}=\frac{16}{25}.\] We add the numerator and denominator of our resulting fraction giving us the answer $41$. $\square$ -mark888

 


	\item (\textit{AIME I 2017}) Let $a_{10} = 10$, and for each positive integer $n >10$ let $a_n = 100a_{n - 1} + n$. Find the least positive $n > 10$ such that $a_n$ is a multiple of $99$.


 
 \textbf{Answer: }045



\textbf{Solution 1}

Writing out the recursive statement for $a_n, a_{n-1}, \dots, a_{10}$ and summing them gives 
 \[a_n+\dots+a_{10}=(a_n+\cdots+a_{11})+a_{10}=100(a_{n-1}+\dots+a_{10})+n+\dots+10\]Which simplifies to 
 \[a_n=99(a_{n-1}+\dots+a_{10})+\frac{1}{2}(n+10)(n-9)\]Therefore, $a_n$ is divisible by 99 if and only if $\frac{1}{2}(n+10)(n-9)$ is divisible by 99, so $(n+10)(n-9)$ needs to be divisible by 9 and 11. Since our goal is just finding the minimum $n$, we can do casework on whether the $n+10$ or $n-9$ is a multiple of 11 (We choose 11 because it is prime). Assume that $n+10$ is a multiple of 11. Writing out a few terms, $n=12, 23, 34, 45$, we see that $n=45$ is the smallest $n$ that works in this case. Next, assume that $n-9$ is a multiple of 11. Writing out a few terms, $n=20, 31, 42, 53$, we see that $n=53$ is the smallest $n$ that works in this case. The smallest $n$ is $\boxed{045}$.

 Note that we can also construct the solution using CRT by assuming either $11$ divides $n+10$ and $9$ divides $n-9$, or $9$ divides $n+10$ and $11$ divides $n-9$, and taking the smaller solution. (We know this because it's impossible for both $n+10$ and $n-9$ to be divisible by 3).

 \textbf{Another way to get the quadratic} Writing out the first couple of terms modulo $99$, we find $a_n = a_{n-1}+n$ so we have $a_{10}=10,a_{11}=21,a_{12}=33,...$ We can compute their finite differences, 
 \[10,21,33,46,...\]
 \[11,12,13,...\]
 \[1,1,...\] Since there are 3 rows we know it is a quadratic and we can continue, by finding the quadratic passing through $(10,10),(11,21),(12,33)$ to get $\frac{(n^2+n-90)}{2}.$-mathkiddus

 



\textbf{Solution 2 (Modular Systems and all solutions)}

Hmmm... let's write our some of the terms, and convert into terms with $a_{10}$, as it is the only value we know.
 \[a_{11} = 100a_{10} + 11\]
 \[a_{12} = 100a_{11} + 12\]
 \[\Rightarrow 100(100a_{10}+11)+12\]
 \[a_{13} = 100a_{12} + 13\]
 \[\Rightarrow 100(100a_{11}+12)+13\]
 \[\Rightarrow 100(100(100a_{10}+11)+12)+13\]Aha! We see it now! Notice that,
 \[a_n = 100^{n-10}a_{10} + \left( \frac{n(n+1)}{2} - \frac{10(11)}{2} \right)\]
 \[\Rightarrow a_n=100^{n-10}a_{10} + \left( \frac{n(n+1)}{2} -55 \right)\]
 \[\Rightarrow a_n = 10^{2(n-10)}a_{10}+\left(\frac{n(n+1)}{2}-55\right)\]
 \[\Rightarrow a_n=(10^{2(n-10)})a_{10}+\left(\frac{n(n+1)}{2}-55\right)\]
 \[\Rightarrow a_n=(10^{2n-19})+\left(\frac{n(n+1)}{2}-55\right)\]

 Yay! We need $a_n \mod 99 \equiv 0 \mod 99 = 0$. Recall $n>10$, so,
 \[\left(10^{2n-19}a_{10}+\left(\frac{n(n+1)}{2}-55\right)\right) \mod 99 = 0\]

 Ugh, more work :-(. Well, it suffices to check mod 9 and mod 11. Thus, we have two cases:

 Case 1: mod 9
 \[\left(10^{2n-19}+\left(\frac{n(n+1)}{2}-55\right)\right) \mod 9\]
 \[10^{2n-19} \mod 9 + \left(\frac{n(n+1)}{2} -55\right) \mod 9\]in which we have an odd 10 power, so an odd 10 power results in
 \[10^{2n-19} \mod 9 \equiv 1 \mod 9 = 1.\]Now,
 \[\frac{n(n+1)}{2} - 55 \mod 9\]
 \[\frac{n(n+1)}{2} \mod 9 - 55 \mod 9\]
 \[\frac{n(n+1)}{2} \mod 9 - 1\]So we have $1-1 + \frac{n(n+1)}{2} \mod 9 = 0$, in which $\frac{n(n+1)}{2} \mod 9 =0 \mod 9$, so $n(n+1) \mod 9 = 0$.

 Checking residues 0 - 8 class mod 9, we have either $n \equiv 0 \mod 9$ or $n \equiv 8 \mod 9$.

 Case 2: mod 11
 \[10^{2n-19} \mod 11 \text{ is apparent when } 10^{2n+1} \mod 11 \text{ so we have } \equiv 10 \mod 11 = 10.\]Then, $\frac{n(n+1)}{2} -55 \mod 11$, in which $\frac{n(n+1)}{2} \mod 11 - 55 \mod 11$, in which $\frac{n(n+1)}{2} \mod 11$.

 We have $10+\frac{n(n+1)}{2} \mod 11 \equiv 0 \mod 11$, so $\frac{n(n+1)}{2} \mod 11 \equiv -10 \mod 11 \equiv 1 \mod 11$

 So, $\frac{n(n+1)}{2} \mod 11 \equiv 1 \mod 11$, in which $n(n+1) \mod 11 \equiv 2 \mod 11$. So, checking class residuals 0-10, we see that either $n \equiv 1 \mod 11$ or $n \equiv 9 \mod 11$.

 Nice! We found our solution sets! Now, using Chinese Remainder Theorem, our solutions are

 $A$: $n \equiv \{0 \mod 9, 1 \mod 11\}$
$B$: $n \equiv \{0 \mod 9, 9 \mod 11\}$
$C$: $n \equiv \{8 \mod 9, 1 \mod 11\}$
$D$: $n \equiv \{8 \mod 9, 9 \mod 11\}$

 In which our respective values $n\in\{A,B,C,D\}$ are $\{45, 9, 89, 53\}$. Because $n > 10$, $\neg 9$, and the smallest $n$ is $\boxed{045}$ :-).

 



\textbf{Solution 3}


 \[a_n \equiv a_{n-1} + n \pmod {99}\]By looking at the first few terms, we can see that 
 \[a_n \equiv 10+11+12+ \dots + n \pmod {99}\]This implies
 \[a_n \equiv \frac{n(n+1)}{2} - \frac{10*9}{2} \pmod {99}\]Since $a_n \equiv 0 \pmod {99}$, we can rewrite the equivalence, and simplify 
 \[0 \equiv \frac{n(n+1)}{2} - \frac{10*9}{2} \pmod {99}\]
 \[0 \equiv n(n+1) - 90 \pmod {99}\]
 \[0 \equiv 4n^2+4n+36 \pmod {99}\]
 \[0 \equiv (2n+1)^2+35 \pmod {99}\]
 \[64 \equiv (2n+1)^2 \pmod {99}\]The only squares that are congruent to $64 \pmod {99}$ are $(\pm 8)^2$ and $(\pm 19)^2$, so 
 \[2n+1 \equiv -8, 8, 19, \text{or } {-19} \pmod {99}\]$2n+1 \equiv -8 \pmod {99}$ yields $n=45$ as the smallest integer solution.

 $2n+1 \equiv 8 \pmod {99}$ yields $n=53$ as the smallest integer solution.

 $2n+1 \equiv -19 \pmod {99}$ yields $n=89$ as the smallest integer solution.

 $2n+1 \equiv 19 \pmod {99}$ yields $n=9$ as the smallest integer solution. However, $n$ must be greater than $10$.

 The smallest positive integer solution greater than $10$ is $n=\boxed{045}$.

 



\textbf{Solution 4 (quadratic)}

Take the mods of the first few terms: $a_{10} \equiv 10 \pmod{99}$, $a_{11} \equiv 21 \pmod{99}$, $a_{12} \equiv 33 \pmod{99}$, $a_{13} \equiv 46 \pmod{99}$, etc.

 Notice the pattern: We're adding $10$ to the mods, then $11$, then $12$, and so on. So we have $a_{n} \equiv n + a_{n - 1} \pmod{99}.$ Using the sum of integers formula, we can create an equation for any general $a_n$ and its mod 99:

 
 \[10 + 11 + \dots + n = \frac{n \cdot (n + 1)}{2} - 45 \implies a_n \equiv \frac{n \cdot (n + 1)}{2} - 45 \pmod{99}.\] Let $\frac{n \cdot (n + 1)}{2} - 45$ equal some multiple of 99 (we'll call this number $99x$). We get the quadratic $n^2 + n - (90 + 198x) = 0.$ Using the quadratic formula we solve for $n$: 
 \[n = \frac{-1 \pm \sqrt{1^2 - 4(1)(-(90+198x))}}{2} = \frac{-1 \pm \sqrt{1 + 360 + 792x}}{2} = \frac{-1 \pm \sqrt{361 + 792x}}{2}.\] So now the problem becomes finding the least $x$ such that $361 + 792x$ is a perfect square, say $p^2.$ Rearranging gives $(p - 19)(p + 19) = 792x.$ We can prime factorize $792$ to get $2^3 \cdot 3^2 \cdot 11.$ So we need to multiply this with some number $x$ so that two factors multiplying to $792x$ differ by $38.$ We find that the least $x$ that works is $x = 10$ (we can have $110 \cdot 72 = 7920$). Then $n = \frac{-1 \pm 91}{2} = \xcancel{-46}, \boxed{45}.$




\textbf{Solution 5}

$a_n=a_{n-1} + n \pmod{99}$.  Using the steps of the previous solution we get up to $n^2+n \equiv 90 \pmod{99}$. This gives away the fact that $(n)(n+1) \equiv 0 \pmod{9} \implies n \equiv \{0, 8\} \pmod{9}$ so either $n$ or $n+1$ must be a multiple of 9. 

 Case 1 ($9 \mid n$): Say $n=9x$ and after simplification $x(9x+1) = 10 \pmod{90} \forall x \in \mathbb{Z}$.  

 Case 2: ($9 \mid n+1$): Say $n=9a-1$ and after simplification $(9a-1)(a) = 10 \pmod{90} \forall a \in \mathbb{Z}$.

 As a result $a$ must be a divisor of $10$ and after doing some testing in both cases the smallest value that works is $x=5 \implies \boxed{045}$.

 



\textbf{Solution 5.5 (not good, risky)}

We just notice that $100 \equiv 1 \pmod{99}$, so we are just trying to find $10 + 11 + 12 + \cdots + n$ modulo $99$, or $\dfrac{n(n+1)}{2} - 45$ modulo $99$. Also, the sum to $44$ is divisible by $99$, and is the first one that is. Thus, if we sum to $45$ the $45$ is cut off and thus is just a sum to $44$.

 Without checking whether there are other sums congruent to $45 \pmod{99}$, we can just write the answer to be $\boxed{045}$.

 



\textbf{Solution 6}

Let $b_n = 2a_{n+10}$. We can find a formula for $b_n$: 

 $b_n = (20+n)(n+1)$. 

 Notice that both can't have a factor of 3. Thus we can limit our search range of n to $n \equiv 7,8 \pmod{9}$. Testing values for n in our search range (like 7,8,16,17,25,26...), we get that 35 is the least n. But, don't write that down! Remember, $b_n = 2a_{n+10}$, so, the 35th term in b corresponds to the 45th term in a. Thus our answer is $\boxed{045}$.

 -AlexLikeMath

 



\textbf{Solution 7 (slower and safer bash)}

The first thing you should realize is that each term after the tenth is another two-digit number chained to the last number. $10, 1011, 101112, \dots$. Now the fact that the sequence starts at $10$ can be completely discarded for this solution. Just consider $a(10)$ then same as $a(1)$, and we can add nine to the answer at the end.

 
The second step is to split $99$ as $9$ and $11$ and solve for divisibility rules individually. Let's start with $11$ because it gives us the most information to continue.

 
In any number generated, if the numbers don't go beyond $20$, then the highest number we can get is $10111213141516171819$, with every odd digit being $1$. This is a little risky because we are assuming that it doesn't exceed $20$. If someone wanted to be absolutely sure they could continue, but this is unnecessary later and a big hassle. Anyways, now we write an equation to check for divisibility by $11$. The expression being $\frac{((n-1) + 0)\cdot n)}{2}-n$.

 The concept here is to add $0$ to the $n-1^{th}$ term altogether, then subtract the number of ones in it, which is $n$. Simplify to $\frac{n(n-3)}{2}$ congruent to $0 \pmod{11}$. Now notice the divide by two can be discarded because one of $n$ or $(n-3)$ will be even. So if $n$ or $n-3$ is to be divisible by $11$, we can make a simple list. 

 
 \[n = 0, 3, 11, 14, 22, 25, 33, 36, 44, 47, \dots\]

 Now we test each $n$ for divisibility by $9$.This is done by making a list that ultimately calculates the sum of every digit in the large number.$n(1)$ to $n(10)$ has the first digit $1$. $n(11)$ to $n(20)$ has the first digit $2$, and so on.The necessary thing to realize is that the sum of all digits $0-9$ is divisible by $9$, so we only have to solve for the sum of the first digits, and then the short list of second digits.

 For example, let's test $n=25$.

 So we know that $25$ include both $1-10$ and $11-20$, so that's $10 + 20$ right away. $21-25$ contains $5$ numbers that have the first digit $3$, so $+15$. Then we add $0-4$ together, which is $10$. $10+20+15+10=55$, which is not divisible by $9$, so it is not the answer. 

 Do this for just a minute you get that $36$ sums to $99$, a multiple of nine! So $n(36)$ is the answer, right? Don't forget we have to add $9$ because we translated $n(10)$ to $n(1)$ at the very beginning!Finally, after a short bash, we get $\boxed{045}$.

 -jackshi2006($LaTeX$ by PureSwag)

 



\textbf{Solution 8 (Timekilling plain Bash!!!!!!!!!!!!!!!!!)}

We will work$\mod 99$. The recursive formula now becomes $a_n=a_{n-1}+n$. Now, we will bash. For convenience, everything is taken $\mod 99$. The sequence is
 \begin{align*}a_{10}&=10 \\ a_{11}&=10+11=21 \\ a_{12}&=12+21=33 \\ a_{13}&=13+33=46 \\ a_{14}&=46+14=60 \\ a_{15}&=60+15=75 \\ a_{16}&=75+16=91 \\ a_{17}&=91+17=9 \\ a_{18}&=9+18=27 \\ a_{19}&=27+19=46 \\ a_{20}&=46+20=66 \\ a_{21}&=66+21=87 \\ a_{22}&=87+22=10 \\ a_{23}&=10+23=33 \\ a_{24}&=33+24=57 \\ a_{25}&=57+25=82 \\ a_{26}&=82+26=9 \\ a_{27}&=9+27=36 \\ a_{28}&=36+28=64 \\ a_{29}&=64+29=93 \\ a_{30}&=93+30=24 \\ a_{31}&=24+31=55 \\ a_{32}&=55+32=87 \\ a_{33}&=87+33=21 \\ a_{34}&=21+34=55 \\ a_{35}&=55+35=90 \\ a_{36}&=90+36=27 \\ a_{37}&=27+37=64 \\ a_{38}&=64+38=3 \\ a_{39}&=3+39=42 \\ a_{40}&=42+40=82 \\ a_{41}&=82+41=24 \\ a_{42}&=24+42=66 \\ a_{43}&=66+43=10 \\ a_{44}&=10+44=54 \\ a_{45}&=54+45=0.\end{align*}
 Hence the least number $n$ is $n=45$.

 



\textbf{Solution 9}

Taking the recurrence mod $99$, we have 
 \[a_n=a_{n-1}+n\] for $a_{10}=10$. Then, we have 
 \[a_n=10+11+12+\cdots+n \implies a_n=\frac{n(n+1)}2-45 \equiv 0 \pmod{99} \implies n(n+1)-90 \equiv 0 \pmod{99} \implies n(n+1)+9 \equiv 0 \pmod{99} \implies n \equiv 0 \pmod{9}.\] Then, we simply can test these values of $n$ to find that $n=\boxed{045}$ produces a value that is also divisible by $11$.

 -A1001

 \textbf{Note} By the way, if you're wondering, $a_{45}$ is the $72$-digit number 
 \[101\text{,}112\text{,}131\text{,}415\text{,}161\text{,}718\text{,}192\text{,}021\text{,}222\text{,}324\text{,}252\text{,}627\text{,}282\text{,}930\text{,}313\text{,}233\text{,}343\text{,}536\text{,}373\text{,}839\text{,}404\text{,}142\text{,}434\text{,}445.\] The prime factorization of $a_{45}$ is 
 \[3^{2}~\cdot~5~\cdot~11~\cdot~151~\cdot~1381~\cdot~1559~\cdot~1\text{,}511\text{,}647~\cdot~448\text{,}966\text{,}261\text{,}198\text{,}213\text{,}862\text{,}368\text{,}469~\cdot~925\text{,}800\text{,}120\text{,}162\text{,}193\text{,}934\text{,}310\text{,}647\text{,}599\text{,}013\] and $\frac{a_{45}}{99}$ is the $70$-digit number 
 \[1\text{,}021\text{,}334\text{,}660\text{,}759\text{,}209\text{,}274\text{,}666\text{,}881\text{,}033\text{,}578\text{,}309\text{,}366\text{,}494\text{,}245\text{,}588\text{,}215\text{,}591\text{,}276\text{,}503\text{,}428\text{,}324\text{,}671\text{,}055.\]

 


 \textbf{Question} Is there an efficient way to notice that the only squares that are congruent to $64 \pmod {99}$ are $(\pm 8)^2$ and $(\pm 19)^2$?  (In Solution 3)

 \textbf{Answer:} Yes.

 We will solve the question by looking for solutions for $m\in[-49,-48,\cdots,-1,0,1,\cdots,48,49]$. Let the square be $m$, thus satisfying $m^2-64\equiv 0\pmod{99}$  and $m^2-64=(m+8)(m-8)=99k$ for some integer $k$. Checking the first factor, $(m+8)$, for factors of $11$ yields:

 $m+8=0\Rightarrow m-8=-16$
$m+8=11\Rightarrow m-8=-5$
$m+8=22\Rightarrow m-8=6$
$m+8=33\Rightarrow m-8=17$
$m+8=44\Rightarrow m-8=28$

 In the first case, the product does divide $99$, so $m=-8\equiv91$ is a solution. For the others, since the first factor already divides $11$, the second factor must divide $9$ (or $3$ in the case of $m+8=33$, which already has a factor of $3$) in order for the product to divide $99$. Here, that is not the case, so $m=-8$ is the only possible solution.

 Now we check the second factor, $(m-8)$:

 $m-8=0\Rightarrow m+8=16$
$m-8=11\Rightarrow m+8=27$
$m-8=22\Rightarrow m+8=38$
$m-8=33\Rightarrow m+8=49$
$m-8=44\Rightarrow m+8=60$

 We immediately see that $m=8$ is a solution. Using a similar argument as above for the others, we notice that $m=19$ is the only solution in this group. Using the property that $ab\equiv(-a)(-b)\mod99$, it is clear that $m=-19$ is also a solution. As a result, $m=\pm8$ and $m=\pm19$ are the only solutions. (This process takes a much shorter time than it seems.)

 


	\item (\textit{AIME I 2017}) Let $z_1 = 18 + 83i$, $z_2 = 18 + 39i,$ and $z_3 = 78 + 99i,$ where $i = \sqrt{-1}$. Let $z$ be the unique complex number with the properties that $\frac{z_3 - z_1}{z_2 - z_1} \cdot \frac{z - z_2}{z - z_3}$ is a real number and the imaginary part of $z$ is the greatest possible. Find the real part of $z$.


 
 \textbf{Answer: }056



\textbf{Solution 1 (Coordinates, Geometry)}

This problem is pretty obvious how to bash, and indeed many of the solutions below explain how to do that. But there’s no fun in that, and let’s see if we can come up with a slicker solution that will be more enjoyable.

 Instead of thinking of complex numbers as purely a real plus a constant times $i$, let’s graph them and hope that the geometric visualization adds insight to the problem.

 
                    \begin{figure}[H]
                        \centering
                        \includegraphics[width=0.5\textwidth]{Solution/AIME_I_2017_10_Solution_Figure_1.png}
                    \end{figure}
                    

 Note that when we subtract two vectors, the geometric result is the line segment between the two endpoints of the vectors. Thus we can fill in $z_3 - z_1, z_2 - z_1, z-z_2$ and $z-z_3$ as so;

 
                    \begin{figure}[H]
                        \centering
                        \includegraphics[width=0.5\textwidth]{Solution/AIME_I_2017_10_Solution_Figure_2.png}
                    \end{figure}
                    
$\angle Z_3Z_1Z_2$ looks similar to $\angle Z_3ZZ_2$, so let’s try to prove that they are congruent. We can show this in two ways;

 



\textbf{Solution 1.1}

Let’s look back at the information given to us. The problem states that $\frac{z_3-z_1}{z_2-z_1}~\cdot~\frac{z-z_2}{z-z_3}$ is a real number. Let the real number be some constant, $c$. Rearranging yields $\frac{z_3-z_1}{z_2-z_1} = c~\cdot~\frac{z-z_3}{z-z_2}$.But how do we relate this expression to our angles? Well, let’s take a look at the divisions themselves.

 The subtraction of two vectors yields a vector, and we can write any vector division as $\frac{V_1}{V_2}= X$ where $X$ is a complex number, as the division of two vectors also yields a vector. We can rewrite this as $V_1 = V_2 \cdot X$. We can think of this expression as transforming $V_2$ directly on to $V_1$, and $X$ is the transformation function. However, this transformation must be some kind of rotation, which means that the degree measure of $X$ is equal to the angle between the two vectors since we need to rotate by that angle to lay $V_2$ flat on $V_1$.

 Thus we can rewrite our previous equation as $X_1 = c \cdot X_2$, where the angle of $X_1$ equals the angle between $z_3-z_1$ and $z_2-z_1$ and likewise for $X_2$. More precisely, we can write $X_1, X_2$ as $r_1 \cdot e^{i \theta_1}$ and $r_2 \cdot e^{i\theta_2}$, respectively by Euler’s formula. Then $\theta_1= \theta_2$ is the claim we wish to prove.

 We can now do some simple algebra to prove this;  

 $r_1 \cdot e^{i \theta_1} = c \cdot r_2 \cdot e^{i\theta_2} \implies \frac{c \cdot r_2}{r_1} = e^{i(\theta_1-\theta_2)} = \cos{(\theta_1-\theta_2)} + i \cdot \sin{(\theta_1-\theta_2)}$
$\frac{c \cdot r_2}{r_1}$ is obviously real, so $\cos{(\theta_1-\theta_2)} + i \cdot \sin{(\theta_1-\theta_2)}$ must be real as well. But the only way that can happen is if $\sin{(\theta_1-\theta_2)} = 0 \implies \theta_1-\theta_2 = 0 \implies \boxed{\theta_1 = \theta_2}$.

 Food For Thought: Let a, b, c, and d be pairwise distinct complex numbers. Then, a, b, c, and d are concyclic if and only if d-b/d-a : c-b/c-a  is a real number. How can we use the above theorem to prove this?

 



\textbf{Solution 1.2}

Let us write $\frac{z_3 - z_1}{z_2 - z_1}$ as some complex number with form $r_1 (\cos \theta_1 + i \sin \theta_1).$ Similarly, we can write $\frac{z-z_2}{z-z_3}$ as some $r_2 (\cos \theta_2 + i \sin \theta_2).$

 The product must be real, so we have that $r_1 r_2 (\cos \theta_1 + i \sin \theta_1) (\cos \theta_2 + i \sin \theta_2)$ is real. $r_1 r_2$ is real by definition, so dividing the real number above by $r_1 r_2$ will still yield a real number. (Note that we can see that $r_1 r_2 \not= 0$ from the definitions of $z_1,$ $z_2,$ and $z_3$). Thus we have 

 
 \[(\cos \theta_1 + i \sin \theta_1) (\cos \theta_2 + i \sin \theta_2) = \cos \theta_1 \cos \theta_2 - \sin \theta_1 \sin \theta_2 + i(\cos \theta_1 \sin \theta_2 + \cos \theta_2 \sin \theta_1)\]

 is real. The imaginary part of this is $(\cos \theta_1 \sin \theta_2 + \cos \theta_2 \sin \theta_1),$ which we recognize as $\sin(\theta_1 + \theta_2).$ This is only $0$ when $\theta_1 + \theta_2 = k\pi$ for some integer $k$. Here $\theta_2$ represents the major angle$\angle Z_3ZZ_2$, and the angle we want is the minor angle, so we can rewrite the equation as $\theta_1 + 2\pi - \theta_2 = k\pi \implies \theta_1 - \theta_2 = (k-2)\pi$. We can see from the diagram that both $\theta_1$ and $\theta_2$ are obtuse, so therefore $\theta_1 - \theta_2 = 0 \implies \boxed{\theta_1 = \theta_2}$.

 



\textbf{Solution 1 Rejoined}

Now that we’ve proven this fact, we know that all four points lie on a circle (intuitively one can also observe this because $z=z_1$ and $z=z_2$ satisfy the property in the question, and $z=z_3$ techincally gives no imaginary part), so let’s draw that in;

 

                    \begin{figure}[H]
                        \centering
                        \includegraphics[width=0.5\textwidth]{Solution/AIME_I_2017_10_Solution_Figure_3.png}
                    \end{figure}
                    

 While $Z_1, Z_2, Z_3$ are fixed, $Z$ can be anywhere on the circle because those are the only values of $Z$ that satisfy the problem requirements. However, we want to find the real part of the $Z$ with the maximum imaginary part. This Z would lie directly above the center of the circle, and thus the real part would be the same as the x-value of the center of the circle. So all we have to do is find this value and we’re done.

 Consider the perpendicular bisectors of $Z_1Z_2$ and $Z_2Z_3$. Since any chord can be perpendicularly bisected by a radius of a circle, these two lines both intersect at the center. Since $Z_1Z_2$ is vertical, the perpendicular bisector will be horizontal and pass through the midpoint of this line, which is (18, 61). Therefore the equation for this line is $y=61$. $Z_2Z_3$ is nice because it turns out the differences in the x and y values are both equal (60) which means that the slope of the line is 1. The slope of the perpendicular bisector is therefore -1 and it passes through the midpoint, (48,69), so the equation of this line is $y=-x+117$. Finally, equating the two yields 
 \[61=-x+117 \implies x = \boxed{056}\]

 \textbf{Bashy Solution :)} We know that
 \[z_3-z_1 = (78+99i)-(18+83i) = 60 + 16i.\]
 \[z_2-z_1=(18+39i)-(18+83i) = -44i.\]Hence,
 \[\frac{z_3-z_1}{z_2-z_1} = \frac{60 + 16i}{-44i} = \frac{15i-4}{11} = \frac{c}{15i+4}\]where $c \in R$.Let $z = a+bi$. Then,
 \[\frac{z-z_2}{z-z_3} = \frac{(a+bi)-(18+39i)}{(a+bi)-(78+99i)} =\]
 \[\frac{z-z_2}{z-z_3} = \frac{(a-18)+i(b-39)}{(a-78)+i(b-99)} =\]
 \[\frac{z-z_2}{z-z_3} = \frac{((a-18)+i(b-39)((a-78)+i(99-b))}{((a-78)+i(b-99))((a-78)+i(99-b)} =\]
 \[\frac{z-z_2}{z-z_3} = \frac{((a-18)+i(b-39)((a-78)+i(99-b))}{k}\]The numerator is:
 \[(a-18)(a-78)-(b-39)(99-b)+i((b-39)(a-78) + (99-b)(a-18))=\]
 \[a^2+b^2-96a-138b+18 \cdot 78 + 39 \cdot 99 + i(60a - 60b + 39 \cdot 78 + 18 \cdot 99)\]The ratio of the imaginary part to the real part must be $\frac{15}{4}$ because $\frac{z_3-z_1}{z_2-z_1} = \frac{c}{15i+4}.$ Hence,
 \[\frac{60a - 60b + 39 \cdot 78 - 18 \cdot 99}{a^2+b^2-96a-138b+18 \cdot 78 + 39 \cdot 99} = \frac{15}{4} \implies\]
 \[\frac{4a-4b+84}{a^2+b^2-96a-138b+9(2 \cdot 78 + 39 \cdot 11)} = \frac{1}{4} \implies\]
 \[16a - 16b + 336 = a^2 + b^2 -96a - 138b +5265 \implies\]
 \[0 = a^2 + b^2 -112a-122b+4929.\]Evidently, $b$ is maximized when $112a-a^2$ is maximized or when $a = \boxed{056}.$

 fixed a single wrong digit - ethanhansummerfun

 



\textbf{Solution 3}

Algebra Bash

 First we calculate $\frac{z_3 - z_1}{z_2 - z_1}$ , which becomes $\frac{15i-4}{11}$.

 Next, we define $z$ to be $a+bi$ for some real numbers $a$ and $b$. Then, $\frac {z-z_2}{z-z_3}$ can be written as $\frac{(a-18)+(b-39)i}{(a-78)+(b-99)i}.$ Multiplying both the numerator and denominator by the conjugate of the denominator, we get:

 $\frac {[(a-18)(a-78)+(b-39)(b-99)]+[(a-78)(b-39)-(a-18)(b-99)]i}{(a-78)^2+(b-99)^2}$

 In order for the product to be a real number, since both denominators are real numbers, we must have the numerator of $\frac {z-z_2}{z-z_3}$ be a multiple of the conjugate of $15i-4$, namely $-15i-4$ So, we have $(a-18)(a-78)+(b-39)(b-99) = -4k$ and $(a-78)(b-39)-(a-18)(b-99) = -15k$ for some real number $k$.

 Then, we get:$(a-18)(a-78)+(b-39)(b-99) = \frac{4}{15}[(a-78)(b-39)-(a-18)(b-99)]$

 Expanding both sides and combining like terms, we get:

 $a^2 - 112a +b^2 - 122b + 4929 = 0$

 which can be rewritten as:

 $(a-56)^2 + (b-61)^2 = 1928$

 Now, common sense tells us that to maximize $b$, we would need to maximize $(b-61)^2$. Therefore, we must set $(a-56)^2$ to its lowest value, namely 0. Therefore, $a$ must be $\boxed{056}.$

 You can also notice that the ab terms cancel out so all you need is the x-coordinate of the center and only expand the a parts of the equation.

 



\textbf{Solution 4 (algebra but much cleaner)}

We see that $\frac{z_3-z_1}{z_2-z_1}=\frac{60+16i}{-44i}=\frac{15i-4}{11}$. 

 Now, let $z-z_3=a+bi$, in which case $z=(a+78)+(b+99)i$ and $z-z_2=(a+60)+(b+60)i$. 

 We now have that $\left(\frac{(a+60)+(b+60)i}{a+bi}\right)\left(\frac{15i-4}{11}\right)$ is real, meaning that:
 \[((a+60)+(b+60)i)\left(\frac{15i-4}{a+bi}\right)\]is also real. 

 We see that:
 \[\frac{15i-4}{a+bi}=\frac{(15i-4)(a-bi)}{a^2+b^2}=\frac{(15b-4a)+(15a+4b)i}{a^2+b^2},\]so therefore, $x=((a+60)+(b+60)i)((15b-4a)+(15a+4b)i)$ is real. 

 This means that $Im(x)=0$, so we now have that:
 \[(a+60)(15a+4b)+(b+60)(15b-4a)=15a^2+15b^2+660a+1140b=0 \Rightarrow a^2+b^2+44a+76b=0.\] 

 This can be rewritten as:
 \[(a+22)^2+(b+38)^2=22^2+38^2.\] In order to maximize $Im(z)$ we want to maximize $b$, and in order to maximize $b$ we want $a+22=0$ and $a=-22$, so $Re(z)=a+78=-22+78=\boxed{056}$. (Note: $Im(\omega)$ is the imaginary part of $\omega$, and $Re(\omega)$ is the real part of $\omega$)~Stormersyle

 



\textbf{Solution 5}

We will just bash. Let $z=a+bi$ where $a,b\in\mathbb{R}$. We see that $\frac{z_3-z_1}{z_2-z_1}=\frac{-4+15i}{11}$ after doing some calculations. We also see that $\frac{[(a-18)+(b-39)i][(a-78)-(b-99)i]}{\text{some real stuff}}.$ We note that $[(a-18)+(b-39)i][(a-78)-(b-99)i]$ is a multiple of $-4-15i$ because the numerator has to be real. Thus, expanding it out, we see that $(a-18)(a-78)+(b-39)(b-99)=-4k \\ (a-78)(b-39)-(a-18)(b-99)=-15k.$ Hence, $(a-18)(a-78)+(b-39)(b-99)=\frac{4}{15}[(a-78)(b-39)-(a-18)(b-99)] \implies a^2-96a+b^2-138b+5265=16a-16b+336 \\ (a-56)^2+(b-61)^2=1928.$ To maximize the imaginary part, $(a-56)^2$ must equal $0$ so hence, $a=\boxed{056}$.

 



\textbf{Solution 6 (Euclidean geometry and coordinates)}

If you've read Evan Chen's book EGMO (Euclidean Geometry in Mathematical Olympiads), you know that $\frac{z_3-z_1}{z_2-z_1}~\cdot~\frac{z-z_2}{z-z_3}$ is a real number is just a fancy way of saying that $z, z_1, z_2, z_3$ are points on a cyclic quadrilateral. The proof of this can be found in EGMO or other sources, and this is the key to the problem.

 And the three given points already define the circle that circumscribes that cyclic quadrilateral. So all we need to do is get the equation for that circle in the Cartesian plane, and then maximize the y-coordinate of $z$. 

 To find the center of the circumcircle, find the intersection of the perpendicular bisectors. One of them is just $y=61$, and the second is $y=-x+117$. So we have the point (56, 61) as the center of the circle. The distance from this point to all three points is $\sqrt{22^2+38^2}$. The square of the radius of this circle is $22^2+38^2=1928$. So the equation for the circle is $(x-56)^2+(y-61)^2=1928$. To maximize the y-coordinate, which does not have to be a positive integer, but $x$ does (per AIME rules given it's the answer), by the trivial inequality $(x-56)^2 \geq 0$, so $\boxed{056}$ is the answer.

 


	\item (\textit{AIME I 2017}) Consider arrangements of the $9$ numbers $1, 2, 3, \dots, 9$ in a $3 \times 3$ array. For each such arrangement, let $a_1$, $a_2$, and $a_3$ be the medians of the numbers in rows $1$, $2$, and $3$ respectively, and let $m$ be the median of $\{a_1, a_2, a_3\}$. Let $Q$ be the number of arrangements for which $m = 5$. Find the remainder when $Q$ is divided by $1000$.


 
 \textbf{Answer: }360



\textbf{Solution 1}

Assume that $5 \in \{a_1, a_2, a_3\}$, $m \neq 5$, and WLOG, $\max{(a_1, a_2, a_3)} = 5$. Then we know that the other two medians in $\{a_1, a_2, a_3\}$ and the smallest number of rows 1, 2, and 3 are all less than 5. But there are only 4 numbers less than 5 in $1, 2, 3, \dots, 9$, a Contradiction. Thus, if $5 \in \{a_1, a_2, a_3\}$, then $m = 5$.

 WLOG, assume $5$ is in the upper left corner. One of the two other values in the top row needs to be below $5$, and the other needs to be above $5$. This can be done in $4\cdot4\cdot2=32$ ways.The other $6$ can be arranged in $6!=720$ ways.Finally, accounting for when $5$ is in every other space, our answer is $32\cdot720\cdot9$, which is $207360$. But we only need the last $3$ digits, so $\boxed{360}$ is our answer.

 \textbf{Alternative Proof} \textbf{Claim:} As long as one median is equal to $5,$ the overall median is equal to $5.$

 \textbf{Proof:} Notice that to obtain a median of $5,$ we need one number greater and one number lesser than $5$ along with $5.$

 Thus, we remain with 3 numbers greater than $5$ and 3 less than $5,$ meaning that the overall median, no matter the distribution, of the remaining two rows must be $5.$

 mathboy282

 



\textbf{Solution 2}

(Complementary Counting with probability)

 Notice that m can only equal 4, 5, or 6, and 4 and 6 are symmetric. 

 WLOG let $m=4$

 1. There is a $\frac{15}{28}$ chance that exactly one of 1, 2, 3 is in the same row with 4.

 There are 3 ways to select which of the smaller numbers will get in the row, and then 5 

 ways to select the number larger than 4.

 $\frac{\dbinom{3}{1}\cdot\dbinom{5}{1}}{\dbinom{8}{2}} = \frac{15}{28}$

 2. There is a $\frac{2}{5}$ chance that the other two smaller numbers end up in the same row.

 There are 2 ways to select the row that the two smaller number are in, and then $\dbinom{3}{2}$ ways 

 to place the smaller numbers in the row.

 $\frac{\dbinom{2}{1}\cdot\dbinom{3}{2}}{\dbinom{6}{2}} = \frac{2}{5}$
$9!(1-2*\frac{15}{28}*\frac{2}{5})=362880*\frac{4}{7}=207\boxed{360}$.

 



\textbf{Solution 3}

We will make sure to multiply by $3!$ in the end to account for all the possible permutation of the rows.

 
WLOG, let $5$ be present in the Row \#$1$.

 
Notice that $5$ MUST be placed with a number lower than it and a number higher than it.

 
This happens in $4\cdot4$ ways. You can permutate Row \#$1$ in $3!$ ways.

 
Now, take a look at Row $2$ and Row $3$.

 
Because there are $6$ numbers to choose from now, you can assign \#'s to Row's \#2&3 in 

 
$\frac{\binom{6}{3}\cdot\binom{3}{3}}{2}$ ways. There are $3!\cdot3!$ ways to permute the numbers in the individual Rows.

 Hence, our answer is $3!(4\cdot4\cdot3!\cdot{10}\cdot{3!}\cdot{3!})=3!(16\cdot60\cdot36)=3!(34560)\implies{207\boxed{360}}$




\textbf{Solution 4}

We see that if one of the medians is 5, then there are two remaining numbers greater than 5 and two less than 5, so it follows that $m=5$. There are 3 ways to choose which row to have 5 in, $4\cdot 4=16$ ways to choose the other two numbers in that row, $3!=6$ ways to arrange the numbers in that row, and $6!=720$ ways for the remaining numbers, for our answer is $3\cdot 16\cdot 6\cdot 720=207\boxed{360}$.-Stormersyle

 



\textbf{Solution 5}

We take the grid, and we do a bunch of stuff with it. First, we sort each row, smallest on the left, largest on the right, then we arrange these 3 rows such that the middle \#s are increasing, from top to bottom. Thus, we get that the cell in the very center of the grid must be 5. (We need to multiply by 1296 at the end)

 Let S mean a number that is < 5, L mean a number > 5.

 
                    \begin{figure}[H]
                        \centering
                        \includegraphics[width=0.5\textwidth]{Solution/AIME_I_2017_11_Solution_Figure_1.jpg}
                    \end{figure}
                    

 In the 2 blank corners, one can be S, and the other one has to be L.

 If the top right corner is S, there are 16 ways, otherwise, there are 144 ways. (This is left as an exercise to the reader).

 Thus, there are 160 * 1296 total configurations, which gets us an answer of $207\boxed{360}$

 -AlexLikeMath

 



\textbf{Solution 6}

We note that if $5$ is a median of one of the rows, then $m=5$. First, focus on the row with $5$ in it. There are $4^2$ ways to choose the other numbers in that row and then $3!$ ways to order it. Now, clearly, there are $6!$ ways to put the other $6$ numbers into the remaining slots so $Q=6!\cdot3!\cdot4^2\cdot3=207360$. Hence, our answer is $\boxed{360}$.

 



\textbf{Solution 7}

Let $A = \{1,2,3,4\}$ and $B = \{6,7,8,9\}$. WLOG, if we let 5 be on the first row, then $m=5$ if and only if $a_1=5$. Thus, the two other numbers that are on row 1 $p, q$ satisfies one of the two following conditions:

 
 \[(1) \quad p\in A \quad \textrm{and} \quad q\in B\]or
 \[(2) \quad p\in B \quad \textrm{and} \quad q\in A\]

 The probability of this happening can be calculated through $\frac{4^2}{{8 \choose 2}} = \frac{4}{7}$. Therefore, 
 \[Q = 9! \cdot \frac{4}{7} = 207360 \equiv \boxed{360} \pmod {1000} .\]




\textbf{Solution  8}

We consider all possible combinations for which the medians lie on the 2nd column in increasing order, and the numbers in each row are strictly increasing. There are nine cases for the remaining two medians $(a_1,a_3)$:

 \textbf{Case 1: }$(a_1,a_3)=(2,6)$Note that $7,8,9$ must all lie in the 3rd column. This yields $2!\cdot 3!=12$ combinations.

 \textbf{Case 2: }$(a_1,a_3)=(2,7)$Note that $1$ must lie on the first cell and $8,9$ must all lie on the 3rd column. If $3$ lie on the 3rd column, there are $2!$ ways to permute $8,9$, while the remaining numbers are fixed.

 If $4$ lie on the 3rd column, there are $2!$ ways to permute $8,9$, while the remaining numbers are fixed.

 If $6$ lie on the 3rd column, there are $2!$ ways to choose where $6$ lie, $2!$ ways to permute $8,9$, $2!$ ways to permute $3,4$. There are hence $12$ combinations in this case.

 \textbf{Case 3: }$(a_1,a_3)=(2,8)$There are $2\cdot2=4$ ways to choose the 2rd row, and there are $2!$ ways to permute the remaining numbers. This yields $2\cdot 4=8$ combinations.

 \textbf{Case 4: }$(a_1,a_3)=(3,6)$Note that $7,8,9$ must all lie in the 3rd column, and there are $2!$ ways to choose where $4$ lie. Thus, there are $3!\cdot2!\cdot2=24$ combinations in this case.

 \textbf{Case 5: }$(a_1,a_3)=(3,7)$Note that $1,2$ must lie on the first cell and $8,9$ must all lie on the 3rd column. If $4$ lies on the 1st column, there are $2\cdot2=4$ ways to permute $1,2,4$ as $4$ cannot lie next to 3. There are $4$ ways to permute $6,8,9$ for a similar reason. There are $4\cdot4=16$ combinations.

 If $4$ lies on the 3rd column, it can only lie next to $3$, resulting in $2$ permutations of $8,9$, the 1st column is permuted similarly. This yields $2\cdot 2=4$ combinations. In total, there are $16+4=20$ combinations in this case.

 \textbf{Case 6: }$(a_1,a_3)=(3,8)$Note that $1,2$ must lie on the 1st column and $9$ is fixed. If $4$ lies on the 1st column, there are $2$ ways to choose where $4$ lies, $2$ ways to permute $1,2$, $2$ ways to permute $6,7$. This yields $8$ combinations in this case.

 If $6$ lies on the 1st column, we see that $6$ is fixed and $4,7$ are fixed, from which there are $2$ ways to permute $1,2$.

 If $7$ lies on the 1st column, we see that $7$ is fixed and $4,6$ are fixed, from which there are $2$ ways to permute $1,2$. In total, there are $8+2+2=12$ combinations in this case.

 \textbf{Case 7: }$(a_1,a_3)=(4,6)$Note that $1,2,3$ must lie on the 1st column and $7,8,9$ must all lie in the 3rd column. This yields $3!\cdot3!=36$ combinations in this case.

 \textbf{Case 8: }$(a_1,a_3)=(4,7)$Note that $1,2,3$ must lie on the 1st column and $6,8,9$ must all lie in the 3rd column. There are $2!$ ways to choose where $6$ lies. This yields $3!\cdot 2\cdot 2!=24$ combinations in this case.

 \textbf{Case 9: }$(a_1,a_3)=(4,8)$Note that $1,2,3$ must lie on the 1st column and $6,7,9$ must all lie in the 3rd column, where $9$ is fixed. There are $2!$ ways to permute where $6,7$ lies. This yields $3!\cdot 2=12$ combinations in this case.

 There are $12+12+8+24+20+12+36+24+12=160$ combinations in total. For each of our combinations, there are $3!$ ways to permute each of the rows. There are also $3!$ ways to decide the order of the rows. Therefore, $Q=1296\cdot 160=207360$.

 The requested answer is $\boxed{360}$.

 


	\item (\textit{AIME I 2017}) Call a set $S$ product-free if there do not exist $a, b, c \in S$ (not necessarily distinct) such that $a b = c$. For example, the empty set and the set $\{16, 20\}$ are product-free, whereas the sets $\{4, 16\}$ and $\{2, 8, 16\}$ are not product-free. Find the number of product-free subsets of the set $\{1, 2, 3, 4, \ldots, 7, 8, 9, 10\}$.


 
 \textbf{Answer: }252



\textbf{Solution 1 (Casework)}

We shall solve this problem by doing casework on the lowest element of the subset. Note that the number $1$ cannot be in the subset because $1*1=1$. Let $S$ be a product-free set. If the lowest element of $S$ is $2$, we consider the set $\{3, 6, 9\}$. We see that 5 of these subsets can be a subset of $S$ ($\{3\}$, $\{6\}$, $\{9\}$, $\{6, 9\}$, and the empty set). Now consider the set $\{5, 10\}$. We see that 3 of these subsets can be a subset of $S$ ($\{5\}$, $\{10\}$, and the empty set). Note that $4$ cannot be an element of $S$, because $2$ is. Now consider the set $\{7, 8\}$. All four of these subsets can be a subset of $S$. So if the smallest element of $S$ is $2$, there are $5*3*4=60$ possible such sets.

 If the smallest element of $S$ is $3$, the only restriction we have is that $9$ is not in $S$. This leaves us $2^6=64$ such sets.

 If the smallest element of $S$ is not $2$ or $3$, then $S$ can be any subset of $\{4, 5, 6, 7, 8, 9, 10\}$, including the empty set. This gives us $2^7=128$ such subsets.

 So our answer is $60+64+128=\boxed{252}$.

 



\textbf{Solution 2 (PIE)}

We will consider the $2^9 = 512$ subsets that do not contain 1. A subset is product-free if and only if it does not contain one of the groups $\{2, 4\}, \{3, 9\}, \{2, 3, 6\},$ or $\{2, 5, 10\}$. There are $2^7$ subsets that contain 2 and 4 since each of the seven elements other than 2 and 4 can either be in the subset or not. Similarly, there are $2^7$ subsets that contain 3 and 9, $2^6$ subsets that contain 2, 3, and 6, and $2^6$ subsets that contain 2, 5, and 10. The number of sets that contain one of the groups is:
 \[2^7 + 2^7 + 2^6 + 2^6 = 384\]For sets that contain two of the groups, we have:
 \[2^5 + 2^5 + 2^5 + 2^5 + 2^4 + 2^4 = 160\]For sets that contain three of the groups, we have:
 \[2^4 + 2^3 + 2^3 + 2^3 = 40\]For sets that contain all of the groups, we have:
 \[2^2 = 4\]By the principle of inclusion and exclusion, the number of product-free subsets is
 \[512 - 384 + 160 - 40 + 4 = \boxed{252}\].

 



\textbf{Solution 3}

Let $X$ be a product-free subset, and note that 1 is not in $x$. We consider four cases: 

 1.) both 2 and 3 are not in $X$. Then there are $2^7=128$ possible subsets for this case. 

 2.) 2 is in $X$, but 3 is not. Then 4 is not in $X$, so there are $2^6=64$ subsets; however, there is a $\frac{1}{4}$ chance that 5 and 10 are both in $X$, so there are $64\cdot \frac{3}{4}=48$ subsets for this case. 

 3.) 2 is not in $X$, but 3 is. Then, 9 is not in $X$, so there are $2^6=64$ subsets for this case. 

 4.) 2 and 3 are both in $X$. Then, 4, 6, and 9 are not in $X$, so there are $2^4=16$ total subsets; however, there is a $\frac{1}{4}$ chance that 5 and 10 are both in $X$, so there are $16\cdot \frac{3}{4}=12$ subsets for this case.

 Hence our answer is $128+48+64+12=\boxed{252}$. -Stormersyle

 



\textbf{Solution 4}

First, ignore 7 as it does not affect any of the other numbers; we can multiply by 2 later (its either in or out).

 Next, we do casework based on whether 2, 3, or 5 are in the set ($2^3 = 8$ quick cases). For each of the cases, the numbers will be of two types --- 

 1. They cannot be in the set because they form a product

 2. Whether we add the number to the set or not, it does not affect the rest of the numbers. These numbers simply add a factor of 2 to the case.

 For example, in the case where all three are present, we can see that $4, 6, 9, 10$ are of type 1. However, $8$ is of type 2. Therefore this case contributes 2. 

 Summing over the 8 cases we get $2 + 4 + 8 + 16 + 16 + 16 + 32 + 32 = 126.$ Multiplying by $2$ because of the $7$ gives $252.$




\textbf{Solution 5}

Note that if any element $s\geq6$ makes an invalid set, it must be $c$ of $ab=c$. Otherwise, $ab\geq12>10$ so no $ab=c$ will suffice. Therefore, whether or not an element $s\geq6$ depends only on the previous elements $1\leq\omega\leq5$.

 First, if a set is product-free, it must never contain an element $\omega=1$ or $1\cdot1=1$.

 We do casework on the subset of $S\in{1,2,3,4,5}=s_1$ to determine $S\in{6,7,8,9,10}=s_2$ (There are only 12 to cover so we just list them out):

 When it is empty, there are clearly no restrictions for $s_2$ so there are $2^5=32$ cases.

 $s_1={2}, s_1={4}, s_1={5}$ has no restrictions $s_2$ so we get $2^5\cdot3=96$ total cases.

 $s_1=3$ only cannot contain 9 so there are $2^4=16$ cases.

 $s_1={2, 5}; s_1={3, 4}; s_1={3, 5}$ each have one restriction in $s_2$ so there are $2^4\cdot3=48$ ways.

 $s_1={2, 3}$ has 2 restrictions (6 and 9) so there are $2^3=8$ ways.

 $s_1={4, 5}$ has no restrictions so there are $2^5=32$ ways.

 $s_1={2, 3, 5}$ has 3 restrictions (6, 9, and 10) so there are $2^2=4$ ways.

 $s_1={3, 4, 5}$ has 1 restriction so there are $2^4=16$ cases.

 Summing the numbers, we get $32+96+16+48+8+32+4+16=\boxed{252}$.

 



\textbf{Solution 6 (simpler casework)}

As before, every set with the property cannot contain $(2, 3, 6), (2, 5, 10), (3, 9)$, or $(2, 4)$.

 \textbf{Case 1}: $S$ contains 2.In this case, seperate the numbers 3, 4, ...10 (excluding 4) into three sets, $[5, 10], [3, 6, 9], [7, 8]$.  The first set has 3 choices, since $S$ can contain 5 or 10 but not both.  The second set then has $2^3-3=5$ valid choices, excluding $(3, 6)$ and $(3, 9); (3, 6, 9)$.  7 and 8 do not appear in any of the prohibited pairs, and therefore have $2^2=4$ choices.  Multiplying $3\cdot 5\cdot 4$, we have 60 options.

 \textbf{Case 2}: $S$ does not contain 2.The only excluded pair possible is $(3, 9)$.  Split the set $[3, 4, 5, 6, 7, 8, 9, 10]$ into $[3, 9]$ and $[4, 5, 6, 7, 8, 10]$.  The first has $2^2-1=3$ choices, excluding $(3, 9)$, and the second has $2^6=64$.  Multiplying, we have $192+60=\boxed{252}$.

 


	\item (\textit{AIME I 2017}) For every $m \geq 2$, let $Q(m)$ be the least positive integer with the following property: For every $n \geq Q(m)$, there is always a perfect cube $k^3$ in the range $n < k^3 \leq mn$. Find the remainder when

 \[\sum_{m = 2}^{2017} Q(m)\]is divided by $1000$.


 
 \textbf{Answer: }059



\textbf{Solution 1}

Lemma 1: The ratio between $k^3$ and $(k+1)^3$ decreases as $k$ increases.

 Lemma 2: If the range $(n,mn]$ includes $y$ cubes, $(p,mp]$ will always contain at least $y-1$ cubes for all $p$ in $[n,+\infty)$.

 If $m=14$, the range $(1,14]$ includes one cube. The range $(2,28]$ includes 2 cubes, which fulfills the Lemma. Since $n=1$ also included a cube, we can assume that $Q(m)=1$ for all $m>14$. Two groups of 1000 are included in the sum modulo 1000. They do not count since $Q(m)=1$ for all of them, therefore 
 \[\sum_{m = 2}^{2017} Q(m) \equiv \sum_{m = 2}^{17} Q(m) \mod 1000\]

 Now that we know this we will find the smallest $n$ that causes $(n,mn]$ to contain two cubes and work backwards (recursion) until there is no cube in $(n,mn]$.

 For $m=2$ there are two cubes in $(n,2n]$ for $n=63$. There are no cubes in $(31,62]$ but there is one in $(32,64]$. Therefore $Q(2)=32$.

 For $m=3$ there are two cubes in $(n,3n]$ for $n=22$. There are no cubes in $(8,24]$ but there is one in $(9,27]$. Therefore $Q(3)=9$.

 For $m$ in $\{4,5,6,7\}$ there are two cubes in $(n,4n]$ for $n=7$. There are no cubes in $(1,4]$ but there is one in $(2,8]$. Therefore $Q(4)=2$, and the same for $Q(5)$, $Q(6)$, and $Q(7)$ for a sum of $8$.

 For all other $m$ there is one cube in $(1,8]$, $(2,16]$, $(3,24]$, and there are two in $(4,32]$. Therefore, since there are 10 values of $m$ in the sum, this part sums to $10$.

 When the partial sums are added, we get $\boxed{059}\hspace{2 mm}QED\hspace{2 mm} \blacksquare$

 This solution is brought to you by \href{https://artofproblemsolving.com/wiki/index.php?title=User:A1b2}{a1b2}

 



\textbf{Solution 2}

We claim that $Q(m) = 1$ when $m \ge 8$. 

 When $m \ge 8$, for every $n \ge Q(m) = 1$, we need to prove there exists an integer $k$, such that $n < k^3 \le m \cdot n$.

 That's because $\sqrt[3]{m \cdot n} - \sqrt[3]{n} \ge 2\sqrt[3]{n} - \sqrt[3]{n} = \sqrt[3]{n} \ge 1$, so k exists between $\sqrt[3]{m \cdot n}$  and $\sqrt[3]{n}$
$\sqrt[3]{n} < k \le \sqrt[3]{m \cdot n}$.

 We can then hand evaluate $Q(m)$ for $m = 2,3,4,5,6,7$, and get $Q(2) = 32$, $Q(3) = 9$, and all the others equal 2.

 There are a total of 2010 integers from 8 to 2017.

 
 \[\sum_{m = 2}^{2017} Q(m) \equiv \sum_{m = 2}^{7} Q(m) + 2010 \equiv 32+9+2+2+2+2+10 = \boxed{059} \mod 1000\]

 -AlexLikeMath

 



\textbf{Solution 3}

Note that the problem is just asking for every $m \geq 2$, find the least $n$ such that there lies a perfect cube in between of $n$ and $mn$ for every $n$ after that minimum including the minimum. We consider $m = 2$ first. Then, we need to find the least $n$ such that for every $n$ after this minimum $n$, including the minimum, there lies a perfect cube in between of $n$ and $2n$. Let's denote the range notation as $(n, 2n)$. Hence, we start with $n = 1$ to get the range of $(1, 2)$. Clearly, no perfect cube lies between these numbers. Then we go to $n = 2$ to get $(2, 4)$ and here also nothing works. We can see that the perfect cubes are $1, 8, 27 64, 125, 216, 343, 512, 729, 1000, ...$. Because $1$ can't lie between two positive integers, we need to find the minimum $n$ such that $8$ is between $n$ and $2n$. Clearly, $n = 4$ will yield this case. We have the range $(4, 8)$. This works. We now list every $n$ starting from $4$ and see if all of these ranges host a perfect cube in between them. We have $(5, 10)$ works, $(6, 12)$ works $(7, 14)$ works, but $(8, 16)$ doesn't work. This is because the next perfect cube has to be $27$. Aha! Now we see the case. If we have a range in the following form $(\frac{a^{3}}{2}, a^{3})$, this will clearly host a perfect cube between them which is $a^{3}$. But as we keep going, we need to put the next perfect cube after $a^{3}$ into the range which is $(a + 1)^{3}$. We know that the next range after $(\frac{a^{3}}{2}, a^{3})$ is just $(a^{3}, 2a^{3})$ because we double everything. Hence, $(a + 1)^{3}$ must be in this range. If the perfect cube in between $a^{3}$ and $2a^{3}$ is of the form $a^{3} + \alpha$, we know that this added constant $\alpha$ is greater than or equal to $(a + 1)^{3} - a^{3} = 3a^{2} + 3a + 1$. Hence, we let $\alpha = a$ to solve $a^{3} \geq 3a^{2} + 3a + 1$ and for integer solutions, we claim the answer is $a \geq 4$. We can prove this by induction.

 Claim: $a^{3} \geq 3a^{2} + 3a + 1$ for integer solutions exists when $a \geq 4$.

 Proof: Let's plug in $a = 4$ for our base case. We get: $64 \geq 48 + 12 + 1 = 61$ which holds true. Hence, the inductive hypothesis is that $(a + 1)^{3} \geq 3(a + 1)^{2} + 3(a + 1) + 1$ for all $a \geq 3$. In fact, we simplify to prove $a^{3} \geq 6a + 6$ for all $a \geq 3$. Now, from our inductive assumption, we need to prove that $3a^{2} + 3a + 1 \geq 6a + 6$. Hence, $3a^{2} - 3a \geq 5$ for all $a \geq 3$. Hence, we need to show $a^{2} - a \geq \frac{5}{3}$ for all $a \geq 3$. But since we're proving this in the integer world, this is equivalent to proving that $a^{2} -a \geq 2$ for all $a \geq 3$. Now we need to show that $a^{2} \geq a + 2$. Dividing by $a$ because $a$ is a positive integer means we have to show $a \geq 1 + \frac{2}{a}$. Now since $a \geq 3$, $\frac{2}{a} < 1$ and so the $RHS$ of the inequality can't even exceed $2$ whereas the $LHS$ of the inequality is already exceeding $2$. Therefore, the inductive hypothesis holds true and we are done.

 What we have just shown now is that $a \geq 4$. Now note that the minimum range we found that worked was $(\frac{a^{3}}{2}, a^{3})$ where $n = \frac{a^{3}}{2}$. We plug in $a = 4$ to get $n = 32$ and hence $Q(2) = 32$. We can now apply the same procedure to the other numbers and finish as the above solutions.

 


	\item (\textit{AIME I 2017}) Let $a > 1$ and $x > 1$ satisfy $\log_a(\log_a(\log_a 2) + \log_a 24 - 128) = 128$ and $\log_a(\log_a x) = 256$. Find the remainder when $x$ is divided by $1000$.


 
 \textbf{Answer: }896



\textbf{Solution 1}

The first condition implies

 
 \[a^{128} = \log_a\log_a 2 + \log_a 24 - 128\]

 \[128+a^{128} = \log_a\log_a 2^{24}\]

 \[a^{a^{128}a^{a^{128}}} = 2^{24}\]

 \[\left(a^{a^{128}}\right)^{\left(a^{a^{128}}\right)} = 2^{24} = 8^8\]

 So $a^{a^{128}} = 8$. 

 Putting each side to the power of $128$:

 
 \[\left(a^{128}\right)^{\left(a^{128}\right)} = 8^{128} = 64^{64},\]

 so $a^{128} = 64 \implies a = 2^{\frac{3}{64}}$. Specifically,

 
 \[\log_a(x) = \frac{\log_2(x)}{\log_2(a)} = \frac{64}{3}\log_2(x)\]

 so we have that

 
 \[256 = \log_a(\log_a(x)) = \frac{64}{3}\log_2\left(\frac{64}{3}\log_2(x)\right)\]

 \[12 = \log_2\left(\frac{64}{3}\log_2(x)\right)\]

 \[2^{12} = \frac{64}{3}\log_2(x)\]

 \[192 = \log_2(x)\]

 \[x = 2^{192}\]

 We only wish to find $x\bmod 1000$. To do this, we note that $x\equiv 0\bmod 8$ and now, by the Chinese Remainder Theorem, wish only to find $x\bmod 125$. By \href{https://artofproblemsolving.com/wiki/index.php?title=Euler\%27s\_Totient\_Theorem}{Euler's Totient Theorem}:

 
 \[2^{\phi(125)} = 2^{100} \equiv 1\bmod 125\]

 so

 
 \[2^{192} \equiv \frac{1}{2^8} \equiv \frac{1}{256} \equiv \frac{1}{6} \bmod 125\]

 so we only need to find the inverse of $6 \bmod 125$. It is easy to realize that $6\cdot 21 = 126 \equiv 1\bmod 125$, so 

 
 \[x\equiv 21\bmod 125, x\equiv 0\bmod 8.\]

 Using \href{https://artofproblemsolving.com/wiki/index.php?title=Chinese\_Remainder\_Theorem}{Chinese Remainder Theorem}, we get that $x\equiv \boxed{896}\bmod 1000$, finishing the solution.

 



\textbf{Solution 2 (Another way to find a)}


 \[\log_a(\log_a(\log_a 2) + \log_a 24 - 128) = 128\]

 \[\implies \log_a(\log_a 2))+\log_a(24)=a^{128}+128\]

 \[\implies \log_a(\log_a 2^{24})=a^{128}+128\]

 \[\implies 2^{24}=a^{a^{(a^{128}+128)}}\]

 Obviously letting $a=2^y$ will simplify a lot and to make the $a^{128}$ term simpler, let $a=2^{\frac{y}{128}}$. Then, 

 
 \[2^{24}=2^{\frac{y}{128} \cdot 2^{\frac{y}{128} \cdot (2^y+128)}}=2^{\frac{y}{128} \cdot 2^{y \cdot (2^{y-7}+1)}}\]

 \[\implies 24=\frac{y}{128} \cdot 2^{y \cdot (2^{y-7}+1)}\]

 \[\implies 3 \cdot 2^{10}=y \cdot 2^{y \cdot (2^{y-7}+1)}\]

 Obviously, $y$ is $3$ times a power of $2$. Testing, we see $y=6$ satisfy the equation so $a=2^{\frac{3}{64}}$. Therefore, $x=2^{192} \equiv \boxed{896} \pmod{1000}$ ~\href{https://artofproblemsolving.com/wiki/index.php?title=Ddk001}{Ddk001}

 \textbf{Alternate solution 1} If you've found $x$ but you don't know that much number theory.

 Note $192 = 3 * 2^6$, so what we can do is take $2^3$ and keep squaring it (mod 1000).

 
 \[2^3 = 8\]
 \[2^6 = 8*8 = 64\]
 \[2^{12} = 64*64 \equiv 96\bmod 1000\]
 \[2^{24} \equiv 96*96 \equiv 216\bmod 1000\]
 \[2^{48} \equiv 216*216 \equiv 656\bmod 1000\]
 \[2^{96} \equiv 656*656 \equiv 336\bmod 1000\]
 \[2^{192} \equiv 336*336 \equiv \boxed{896}\bmod 1000\]

 \textbf{Alternate solution 2} Another way to find $x \bmod 1000$ using modular arithmetic.In the same way as solution $1$, we can find that.
 \[x\equiv 21\bmod 125, x\equiv 0\bmod 8.\]
 \[x = 8m = 125n+21\] For some positive integers $m$ and $n$.Taking the equation mod $8$ gives 
 \[5n+5 \equiv 0\bmod 8\]
 \[n \equiv 7\bmod 8\]
 \[n = 8k-1\] For some positive integer $k$.Plug this back into the original equation.
 \[8m = 125(8k-1)+21\]
 \[8m = 1000k-104\]
 \[x = 8m = 1000k - 104\]
 \[x \equiv -104 \equiv 896\bmod 1000\]
 \[x \equiv 896\bmod 1000\]



	\item (\textit{AIME I 2017}) The area of the smallest equilateral triangle with one vertex on each of the sides of the right triangle with side lengths $2\sqrt3$, $5$, and $\sqrt{37}$, as shown, is $\frac{m\sqrt{p}}{n}$, where $m$, $n$, and $p$ are positive integers, $m$ and $n$ are relatively prime, and $p$ is not divisible by the square of any prime. Find $m+n+p$.


 
                    \begin{figure}[H]
                        \centering
                        \includegraphics[width=0.5\textwidth]{Problem/AIME_I_2017_15_Problem_Figure_1.png}
                    \end{figure}
                    

 \textbf{Answer: }145



\textbf{Solution 1}

Lemma: If $x,y$ satisfy $px+qy=1$, then the minimal value of $\sqrt{x^2+y^2}$ is $\frac{1}{\sqrt{p^2+q^2}}$.

 Proof: Recall that the distance between the point $(x_0,y_0)$ and the line $px+qy+r = 0$ is given by $\frac{|px_0+qy_0+r|}{\sqrt{p^2+q^2}}$. In particular, the distance between the origin and any point $(x,y)$ on the line $px+qy=1$ is at least $\frac{1}{\sqrt{p^2+q^2}}$.

 ---

 Let the vertices of the right triangle be $(0,0),(5,0),(0,2\sqrt{3}),$ and let $(a,0),(0,b)$ be the two vertices of the equilateral triangle on the legs of the right triangle. Then, the third vertex of the equilateral triangle is $\left(\frac{a+b\sqrt{3}}{2},\frac{a\sqrt{3}+b}{2}\right)$. This point must lie on the hypotenuse $\frac{x}{5} + \frac{y}{2\sqrt{3}} = 1$, i.e. $a,b$ must satisfy
 \[\frac{a+b\sqrt{3}}{10}+\frac{a\sqrt{3}+b}{4\sqrt{3}} = 1,\]which can be simplified to
 \[\frac{7}{20}a + \frac{11\sqrt{3}}{60}b = 1.\]

 By the lemma, the minimal value of $\sqrt{a^2+b^2}$ is
 \[\frac{1}{\sqrt{\left(\frac{7}{20}\right)^2 + \left(\frac{11\sqrt{3}}{60}\right)^2}} = \frac{10\sqrt{3}}{\sqrt{67}},\]so the minimal area of the equilateral triangle is
 \[\frac{\sqrt{3}}{4} \cdot \left(\frac{10\sqrt{3}}{\sqrt{67}}\right)^2 = \frac{\sqrt{3}}{4} \cdot \frac{300}{67} = \frac{75\sqrt{3}}{67},\]and hence the answer is $75+3+67=\boxed{145}$.

 



\textbf{Solution 2}

Let $AB=2\sqrt{3}, BC=5$, $D$ lies on $BC$, $F$ lies on $AB$ and $E$ lies on $AC$

 Set $D$ as the origin, $BD=a,BF=b$, $F$ can be expressed as $-a+bi$ in argand plane, the distance of $CD$ is $5-a$

 We know that $(-a+bi)\cdot[\cos(-\frac{\pi}{3})+ i\sin(-\frac{\pi}{3})]=(-\frac{a+\sqrt{3}b}{2}+\frac{\sqrt{3}a+b}{2}i)$. We know that the slope of $AC$ is $-\frac{2\sqrt{3}}{5}$, we have that $\frac{5-a-(-\frac{a+\sqrt{3}b}{2})}{(\sqrt{3}a+b)/2}=\frac{5}{2\sqrt{3}}$, after computation, we have $11b+7\sqrt{3}a=20\sqrt{3}$

 Now the rest is easy with C-S inequality, $(a^2+b^2)(147+121)\geq (7\sqrt{3}a+11b)^2, a^2+b^2\geq \frac{300}{67}$ so the smallest area is $\frac{\sqrt{3}}{4}\cdot \frac{300}{67}=\frac{75\sqrt{3}}{67}$, and the answer is $\boxed{145}$




\textbf{Solution 3}

Let $\triangle ABC$ be the right triangle with sides $AB = x$, $AC = y$, and $BC = z$ and right angle at $A$.

 Let an equilateral triangle touch $AB$, $AC$, and $BC$ at $D$, $E$, and $F$ respectively, having side lengths of $c$.

 Now, call $AD$ as $a$ and $AE$ as $b$. Thus, $DB = x-a$ and $EC = y-b$.

 By Law of Sines on triangles $\triangle DBF$ and $ECF$,

 $BF = \frac{z(a\sqrt{3}+b)} {2y}$ and $CF = \frac{z(a+b\sqrt{3})} {2x}$.

 Summing, 

 $BF+CF =  \frac{z(a\sqrt{3}+b)} {2y} + \frac{z(a+b\sqrt{3})} {2x} = BC = z$.

 Now substituting $AB = x = 2\sqrt{3}$, $AC = y = 5$, and $BC = \sqrt{37}$ and solving,$\frac{7a}{20} + \frac{11b\sqrt{3}}{60} = 1$.

 We seek to minimize $[DEF] = c^2 \frac{\sqrt{3}}{4} = (a^2 + b^2) \frac{\sqrt{3}}{4}$.

 This is equivalent to minimizing $a^2+b^2$.

 Using the lemma from solution 1, we conclude that $\sqrt{a^2+b^2} = \frac{10\sqrt{3}}{\sqrt{67}}$

 Thus, $[DEF] = \frac{75\sqrt{3}}{67}$ and our final answer is $\boxed{145}$

 - Awsomness2000

 



\textbf{Solution 4 (Trigonometry)}


                    \begin{figure}[H]
                        \centering
                        \includegraphics[width=0.5\textwidth]{Solution/AIME_I_2017_15_Solution_Figure_1.png}
                    \end{figure}
                    Let $ED = DF = x, AC = 5, \angle BAC = \alpha,  \angle CDE = \beta.$

 Then $\cot \alpha = \frac {5}{2\sqrt{3}},\angle CDF = \beta - 60^\circ,$
 \[\angle AED = \beta - \alpha, CD = x \cos(\beta - 60^\circ).\] 

 By Law of Sines on triangle $\triangle ADE$ we get 
 \[AD = x \frac {\sin (\beta – \alpha)}{\sin{\alpha}}.\]
 \[AD + CD = AC \implies\]
 \[x(\frac {\sin \beta}{\tan{\alpha}} - \cos \beta) +x (\frac {\cos\beta}{2} +\frac{\sin\beta \sqrt{3}}{2}) = AC\]
 \[\implies \frac{AC}{x} = \sin \beta (\cot \alpha + \frac{\sqrt{3}}{2})- \frac{\cos\beta}{2}\] 
 \[a \sin \beta + b \cos \beta \le \sqrt {a^2+b^2} \implies\]
 \[\frac{AC}{x}  \le \sqrt {\left(\cot \alpha + \frac {\sqrt3}{2}\right)^2  + \left(\frac {1}{2}\right)^2} = \sqrt {\cot^2 \alpha + \sqrt {3} \cot \alpha + 1}\]The smallest area $[DEF]$ is 
 \[\frac {\sqrt {3}}{4} \cdot x_{min}^2 = \frac {\sqrt {3} AC^2}{4 (\cot ^2 \alpha + \sqrt {3} \cot \alpha + 1)} = \frac{75 \sqrt{3}}{67}.\]\textbf{vladimir.shelomovskii@gmail.com, vvsss}

 \textbf{Note} $a \sin \beta + b \cos \beta \le \sqrt {a^2+b^2}$ follows from Cauchy-Schwarz. 
 \[a \sin \beta + b \cos \beta \le \sqrt{a^2+b^2}\sqrt{\sin^2+\cos^2}=\sqrt{a^2+b^2}\]~mathboy282

 



\textbf{Solution 5 (Complex numbers)}

We will use complex numbers. Set the vertex at the right angle to be the origin, and set the axes so the other two vertices are $5$ and $2\sqrt{3}i$, respectively. Now let the vertex of the equilateral triangle on the real axis be $a$ and let the vertex of the equilateral triangle on the imaginary axis be $bi$. Then, the third vertex of the equilateral triangle is given by:
 \[(bi-a)e^{-\frac{\pi}{3}i}+a=(bi-a)(\frac{1}{2}-\frac{\sqrt{3}}{2}i)+a=(\frac{a}{2}+\frac{b\sqrt{3}}{2})+(\frac{a\sqrt{3}}{2}+\frac{b}{2})i\].

 For this to be on the hypotenuse of the right triangle, we also have the following:
 \[\frac{\frac{a\sqrt{3}}{2}+\frac{1}{2}}{\frac{a}{2}+\frac{b\sqrt{3}}{2}-5}=-\frac{2\sqrt{3}}{5}\iff 7\sqrt{3}a+11b=20\sqrt{3}\]

 Note that the area of the equilateral triangle is given by $\frac{\sqrt{3}(a^2+b^2)}{4}$, so we seek to minimize $a^2+b^2$. This can be done by using the Cauchy Schwarz Inequality on the relation we derived above:
 \[1200=(7\sqrt{3}a+11b)^2\leq ((7\sqrt{3})^2+11^2)(a^2+b^2)\implies a^2+b^2\geq \frac{1200}{268}\]

 Thus, the minimum we seek is simply $\frac{\sqrt{3}}{4}\cdot\frac{1200}{268}=\frac{75\sqrt{3}}{67}$, so the desired answer is $\boxed{145}$.

 



\textbf{Solution 6}

In the complex plane, let the vertices of the triangle be $a = 5,$ $b = 2i \sqrt{3},$ and $c = 0.$ Let $e$ be one of the vertices, where $e$ is real. A point on the line passing through $a = 5$ and $b = 2i \sqrt{3}$ can be expressed in the form
 \[f = (1 - t) a + tb = 5(1 - t) + 2ti \sqrt{3}.\]We want the third vertex $d$ to lie on the line through $b$ and $c,$ which is the imaginary axis, so its real part is 0.Since the small triangle is equilateral, $d - e = \cos 60^\circ \cdot (f - e),$ or
 \[d - e = \frac{1 + i \sqrt{3}}{2} \cdot (5(1 - t) - e + 2ti \sqrt{3}).\]Then the real part of $d$ is
 \[\frac{5(1 - t) - e}{2} - 3t + e = 0.\]Solving for $t$ in terms of $e,$ we find
 \[t = \frac{e + 5}{11}.\]Then
 \[f = \frac{5(6 - e)}{11} + \frac{2(e + 5) \sqrt{3}}{11} i,\]so
 \[f - e = \frac{30 - 16e}{11} + \frac{2(e + 5) \sqrt{3}}{11} i,\]so
 \begin{align*} |f - e|^2 &= \left( \frac{30 - 16e}{11} \right)^2 + \left( \frac{2(e + 5) \sqrt{3}}{11} \right)^2 \\ &= \frac{268e^2 - 840e + 1200}{121}. \end{align*}
This quadratic is minimized when $e = \frac{840}{2 \cdot 268} = \frac{105}{67},$ and the minimum is $\frac{300}{67},$ so the smallest area of the equilateral triangle is
 \[\frac{\sqrt{3}}{4} \cdot \frac{300}{67} = \boxed{\frac{75 \sqrt{3}}{67}}.\]




\textbf{Solution 7}

We can use complex numbers. Set the origin at the right angle. Let the point on the real axis be $a$ and the point on the imaginary axis be $bi$. Then, we see that $(a-bi)\left(\text{cis}\frac{\pi}{3}\right)+bi=(a-bi)\left(\frac{1}{2}+i\frac{\sqrt{3}}{2}\right)+bi=\left(\frac{1}{2}a+\frac{\sqrt{3}}{2}b\right)+i\left(\frac{\sqrt{3}}{2}a+\frac{1}{2}b\right).$ Now we switch back to Cartesian coordinates. The equation of the hypotenuse is $y=-\frac{2\sqrt{3}}{5}x+2\sqrt{3}.$ This means that the point $\left(\frac{1}{2}a+\frac{\sqrt{3}}{2}b,\frac{\sqrt{3}}{2}a+\frac{1}{2}b\right)$ is on the line. Plugging the numbers in, we have $\frac{\sqrt{3}}{2}a+\frac{1}{2}b=-\frac{\sqrt{3}}{5}a-\frac{3}{5}b+2\sqrt{3} \implies 7\sqrt{3}a+11b=20\sqrt{3}.$ Now, we note that the side length of the equilateral triangle is $a^2+b^2$ so it suffices to minimize that. By Cauchy-Schwarz, we have $(a^2+b^2)(147+121)\geq(7\sqrt{3}a+11b)^2 \implies (a^2+b^2)\geq\frac{300}{67}.$ Thus, the area of the smallest triangle is $\frac{300}{67}\cdot\frac{\sqrt{3}}{4}=\frac{75\sqrt{3}}{67}$ so our desired answer is $\boxed{145}$.

 (Solution by Pleaseletmetwin, but not added to the Wiki by Pleaseletmetwin)

 



\textbf{Solution 8}

Employ the same complex bash as in Solution 5, but instead note that minimizing $x^2+y^2$ is the same as minimizing the distance from 0,0 to x,y, since they are the same quantity. We use point to plane instead, which gives you the required distance.

 



\textbf{Solution 9 (Non Analytic)}

Let $S$ be the triangle with side lengths $2\sqrt{3},~5,$ and $\sqrt{37}$.

 We will think about this problem backwards, by constructing a triangle as large as possible (We will call it $T$, for convenience) which is similar to $S$ with vertices outside of a unit equilateral triangle $\triangle ABC$, such that each vertex of the equilateral triangle lies on a side of $T$. After we find the side lengths of $T$, we will use ratios to trace back towards the original problem.

 First of all, let $\theta = 90^{\circ}$, $\alpha = \arctan\left(\frac{2\sqrt{3}}{5}\right)$, and $\beta = \arctan\left(\frac{5}{2\sqrt{3}}\right)$ (These three angles are simply the angles of triangle $S$; out of these three angles, $\alpha$ is the smallest angle, and $\theta$ is the largest angle). Then let us consider a point $P$ inside $\triangle ABC$ such that $\angle APB = 180^{\circ} - \theta$, $\angle BPC = 180^{\circ} - \alpha$, and $\angle APC = 180^{\circ} - \beta$. Construct the circumcircles $\omega_{AB}, ~\omega_{BC},$ and $\omega_{AC}$ of triangles $APB, ~BPC,$ and $APC$ respectively. 

 From here, we will prove the lemma that if we choose points $X$, $Y$, and $Z$ on circumcircles $\omega_{AB}, ~\omega_{BC},$ and $\omega_{AC}$ respectively such that $X$, $B$, and $Y$ are collinear and $Y$, $C$, and $Z$ are collinear, then $Z$, $A$, and $X$ must be collinear. First of all, if we let $\angle PAX = m$, then $\angle PBX = 180^{\circ} - m$ (by the properties of cyclic quadrilaterals), $\angle PBY = m$ (by adjacent angles), $\angle PCY = 180^{\circ} - m$ (by cyclic quadrilaterals), $\angle PCZ = m$ (adjacent angles), and $\angle PAZ = 180^{\circ} - m$ (cyclic quadrilaterals). Since $\angle PAX$ and $\angle PAZ$ are supplementary, $Z$, $A$, and $X$ are collinear as desired. Hence, $\triangle XYZ$ has an inscribed equilateral triangle $ABC$.

 In addition, now we know that all triangles $XYZ$ (as described above) must be similar to triangle $S$, as $\angle AXB = \theta$ and $\angle BYC = \alpha$, so we have developed $AA$ similarity between the two triangles. Thus, $\triangle XYZ$ is the triangle similar to $S$ which we were desiring. Our goal now is to maximize the length of $XY$, in order to maximize the area of $XYZ$, to achieve our original goal.

 Note that, all triangles $PYX$ are similar to each other if $Y$, $B$, and $X$ are collinear. This is because $\angle PYB$ is constant, and $\angle PXB$ is also a constant value. Then we have $AA$ similarity between this set of triangles. To maximize $XY$, we can instead maximize $PY$, which is simply the diameter of $\omega_{BC}$. From there, we can determine that $\angle PBY = 90^{\circ}$, and with similar logic, $PA$, $PB$, and $PC$ are perpendicular to $ZX$, $XY$, and $YZ$ respectively We have found our desired largest possible triangle $T$.

 All we have to do now is to calculate $YZ$, and use ratios from similar triangles to determine the side length of the equilateral triangle inscribed within $S$. First of all, we will prove that $\angle ZPY = \angle ACB + \angle AXB$. By the properties of cyclic quadrilaterals, $\angle AXB = \angle PAB + \angle PBA$, which means that $\angle ACB + \angle AXB = 180^{\circ} - \angle PAC - \angle PBC$. Now we will show that $\angle ZPY =  180^{\circ} - \angle PAC - \angle PBC$. Note that, by cyclic quadrilaterals, $\angle YZP = \angle PAC$ and $\angle ZYP = \angle PBC$. Hence, $\angle ZPY = 180^{\circ} - \angle PAC - \angle PBC$ (since $\angle ZPY + \angle YZP + \angle ZYP = 180^{\circ}$), proving the aforementioned claim. Then, since $\angle ACB = 60^{\circ}$ and $\angle AXB = \theta = 90^{\circ}$, $\angle ZPY = 150^{\circ}$.

 Now we calculate $PY$ and $PZ$, which are simply the diameters of circumcircles $\omega_{BC}$ and $\omega_{AC}$, respectively. By the extended law of sines, $PY = \frac{BC}{\sin{BPC}} = \frac{\sqrt{37}}{2\sqrt{3}}$ and $PZ = \frac{CA}{\sin{CPA}} = \frac{\sqrt{37}}{5}$.

 We can now solve for $ZY$ with the law of cosines:

 
 \[(ZY)^2 = \frac{37}{25} + \frac{37}{12} - \left(\frac{37}{5\sqrt{3}}\right)\left(-\frac{\sqrt{3}}{2}\right)\]

 \[(ZY)^2 = \frac{37}{25} + \frac{37}{12} + \frac{37}{10}\]

 \[(ZY)^2 = \frac{37 \cdot 67}{300}\]

 \[ZY = \sqrt{37} \cdot \frac{\sqrt{67}}{10\sqrt{3}}\]

 Now we will apply this discovery towards our original triangle $S$. Since the ratio between $ZY$ and the hypotenuse of $S$ is $\frac{\sqrt{67}}{10\sqrt{3}}$, the side length of the equilateral triangle inscribed within $S$ must be $\frac{10\sqrt{3}}{\sqrt{67}}$ (as $S$ is simply as scaled version of $XYZ$, and thus their corresponding inscribed equilateral triangles must be scaled by the same factor). Then the area of the equilateral triangle inscribed within $S$ is $\frac{75\sqrt{3}}{67}$, implying that the answer is $\boxed{145}$.

 \textbf{-Solution by TheBoomBox77}

 



\textbf{Solution 10}

Let the right triangle's lower-left point be at $O(0,0)$. Notice the 2 other points will determine a unique equilateral triangle. Let 2 points be on the $x$-axis ($B$) and the $y$-axis ($A$) and label them $(b, 0)$ and $(0, a)$ respectively.  The third point ($C$) will then be located on the hypotenuse.  We proceed to find the third point's coordinates in terms of $a$ and $b$. 

 
1. Find the slope of $AB$ and take the negative reciprocal of it to find the slope of the line containing $C$.  Notice the line contains the midpoint of $AB$ so we can then have an equation of the line.

 2. Let $AB=x.$  For $ABC$ to be an equilateral triangle, the altitude from $C$ to $AB$ must be $\frac{x\sqrt{3}}{2}.$

 We then have two equations and two variables, so we can solve for $C$'s coordinates. 

 We can find $C(\frac{a+b\sqrt{3}}{2}), (\frac{b+a\sqrt{3}}{2}).$Also, note that $C$ must be on the hypotenuse of the triangle $\frac{x}{5}+\frac{y}{2\sqrt{3}}=1.$  We can plug in $x$ and $y$ as the coordinates of $C$,  which simplifies to 

 
 \[11b+7\sqrt{3}a=20\sqrt{3}.\] 

 We aim to minimize the side length of the triangle, which is $\sqrt{a^2+b^2}.$ Applying the Cauchy inequality gives us 

 
 \[(a^2+b^2)(7\sqrt{3}^2+11^2)\geq (11b+7\sqrt3a)^2 = 1200\] 

 From which we obtain $\sqrt{a^2+b^2} \geq \sqrt{\frac{300}{67}}.$ Thus, the area of the triangle = $\frac{75\sqrt{3}}{67}$ which leads to the answer $75+3+67=\boxed{145}.$

 -hi\_im\_bob

 



\textbf{Solution 11}


                    \begin{figure}[H]
                        \centering
                        \includegraphics[width=0.5\textwidth]{Solution/AIME_I_2017_15_Solution_Figure_2.png}
                    \end{figure}
                    The general solution to the minimal area is as following:

 
 \[A_{min}={{\sqrt{3}m^2n^2}\over{4(m^2+n^2+\sqrt{3}mn)}},\]

 where $m$ and $n$ are the two legs of the right triangle.  In this particular case $m=5$ and $n=2\sqrt{3}$.  When we plug in these two values, we recover the correct answer of $\frac{75\sqrt3}{67}$.

 The contour of the minimal area $A_{min}$ is plotted as a function of leg lengths $m$ and $n$, as shown on the right hand side.

 \textbf{Note} The proof of the formula can be done in a similar fashion as Solution 4, where instead of using specific values, we define variables, $m$ and $n$ in this case, as per the formula.

 



\textbf{Solution 12 (Geometry)}


                    \begin{figure}[H]
                        \centering
                        \includegraphics[width=0.5\textwidth]{Solution/AIME_I_2017_15_Solution_Figure_3.png}
                    \end{figure}
                    Let $ED = DF = x, AC = 5, \angle BAC = \alpha.$Then $\cot \alpha = \frac {5}{2\sqrt{3}}.$
$P -$ midpoint $DE, M -$ midpoint $FE, Q -$ circumcenter $\triangle AE.$
 \[\angle EQM = \alpha, \angle MQE = 90^\circ - \alpha,\]
 \[\angle PEQ =  \angle MQE + 60^\circ =  150^\circ - \alpha.\]
 \[PE = \frac {x}{2}, QE = \frac {x}{2 \sin \alpha} \implies\]
 \[PQ^2 = PE^2 + QE^2 - 2 PE \cdot QE \cdot \cos(150^\circ - \alpha),\]
 \[PQ^2 =  \frac{x^2}{4} (\cot ^2 \alpha + \sqrt {3} \cot \alpha + 1) .\]Points $P$ and $Q$ lies on bisectors of $CE$ and $AE,$ so $PQ \ge \frac {AC}{2}.$

 The smallest area $[DEF]$ is 
 \[\frac {\sqrt {3}}{4} \cdot x_{min}^2 = \frac {\sqrt {3} AC^2}{4 (\cot ^2 \alpha + \sqrt {3} \cot \alpha + 1)} = \frac{75 \sqrt{3}}{67}.\]\textbf{vladimir.shelomovskii@gmail.com, vvsss}

 



\textbf{Solution 13 (Kinematics+Geometry)}


                    \begin{figure}[H]
                        \centering
                        \includegraphics[width=0.5\textwidth]{Solution/AIME_I_2017_15_Solution_Figure_4.png}
                    \end{figure}
                    Let $ED = DF = x, AC = 5, \angle BAC = \alpha.$Then $\cot \alpha = \frac {5}{2\sqrt{3}}.$

 Let the required triangle with minimal sides $DEF$ be constructed.

 Let us vary its position in an acceptable way, that is, we will perform a movement in which the vertices of the $\triangle DEF$ remain on the sides of the $\triangle ABC.$ Any motion of a solid plane figure can be considered as rotation around some point, the center of rotation. The speed of movement of any point is perpendicular to the segment from the center of rotation to this point. With an acceptable variation, the velocities of the vertices of the $\triangle DEF$ are directed along the sides $\triangle ABC.$ Consequently, there is a point $X$ located at the intersection of perpendiculars to the sides of $\triangle ABC,$ around which $\triangle DEF$ rotates. The bases of the perpendiculars dropped from $X$ to the sides of $\triangle ABC$ form a regular $\triangle DEF.$

 Therefore $X$ is the first isodynamic point of $\triangle ABC.$

 It is known that $\angle AXB = \angle ACB + 60^\circ = 150^\circ.$

 The points $A, F, X,$ and $E$ are concyclic, $AX$ is the diameter, so $AX = \frac {x}{\sin \alpha}.$ Similarly,  $BX = \frac {x}{\cos \alpha}.$
 \[AX^2 + BX^2 - 2 AX BX \cos 150^\circ = AB^2, AC = AB \cos \alpha \implies\]
 \[AC^2 = x^2 \cdot  (\cot ^2 \alpha + \sqrt {3} \cot \alpha + 1).\] The smallest area $[DEF]$ is 
 \[\frac {\sqrt {3}}{4} \cdot x_{min}^2 = \frac {\sqrt {3} AC^2}{4 (\cot ^2 \alpha + \sqrt {3} \cot \alpha + 1)} = \frac{75 \sqrt{3}}{67}.\]\textbf{vladimir.shelomovskii@gmail.com, vvsss}

 



\textbf{Solution 14 (Fastest Trig)}

Let the right triangle be $\triangle ABC$ with $BC=2\sqrt{3}$, $AC=5$, and $AB=\sqrt{37}$, and let the equilateral be $\triangle DEF$ with $D, E, F$ on $AC, AB, BC$ respectively. Call its side length $x$. Let the foot from $E$ to $BC$ be $H$, and let $\angle EFH=\theta$. Since $\angle EFD=60^{\circ}$, $\angle DFC=120^{\circ}-\theta$. We get $FC=x\cos(120^{\circ}-\theta)$ and $HF=x\cos\theta$. Meanwhile, since $EH =x\sin\theta$ and $\triangle EBC \sim \triangle ABC$, $BH=\frac{2\sqrt{3}}{5} x\sin\theta$. We have

 
 \[2\sqrt{3}=BC=BH+HF+FC=\frac{2\sqrt{3}}{5}x\sin\theta + x\cos\theta +x\cos\left(120^{\circ}-\theta\right)\]

 \[\implies 2\sqrt{3}=x\left(\frac{2\sqrt{3}}{5}\sin\theta+\cos\theta-\frac{1}{2}\cos\theta+\frac{\sqrt{3}}{2}\sin\theta\right)\]

 \[\implies x=\frac{2\sqrt{3}}{\frac{9\sqrt{3}}{10}\sin\theta+\frac{1}{2}\cos\theta}.\]

 \[\implies [\triangle DEF]=\frac{\sqrt{3}}{4}x^2=\frac{\sqrt{3}}{4}\cdot \frac{12}{\left(\frac{9\sqrt{3}}{10}\sin\theta+\frac{1}{2}\cos\theta\right)^2}=\frac{3\sqrt{3}}{\left(\frac{9\sqrt{3}}{10}\sin\theta+\frac{1}{2}\cos\theta \right)^2}\]

 By Cauchy-Schwarz, 
 \[\left(\frac{9\sqrt{3}}{10}\sin\theta+\frac{1}{2}\cos\theta \right)^2 \le \left(\cos^2\theta+\sin^2\theta\right)\left(\left(\frac{9\sqrt{3}}{10}\right)^2+\left(\frac{1}{2}\right)^2\right)=\frac{67}{25}\]

 \[\implies [\triangle DEF] \ge \frac{3\sqrt{3}}{\frac{67}{25}}=\frac{75\sqrt{3}}{67}.\]

 The requested sum is $75+3+67=\boxed{145}$.

 



\textbf{Solution 15 (Trig & Basic Calc)}

We label the points such that $AB=2\sqrt3$, $BC=5$, and $AC=\sqrt{37}$. Let the point where the equilateral triangle hits $AB$ be $D$, the point where the triangle hits $BC$ be $E$, and the point where the triangle hits $AC$ be $F$. Finally, let the side length of the equilateral triangle be $x$ and $\angle{BDE}=\theta$.

 Note that $\angle{ADF}=120^\circ-\theta$. Simple right triangle trigonometry gives $\sin\angle{A}=\dfrac{5}{\sqrt{37}}$, so we can use Law of Sines on $\triangle{ADF}$ to find that
 \[\dfrac{x\sqrt{37}}{5}=\dfrac{k}{\sin(120^\circ-\theta)},\]where $k=AF$. We can rearrange to get $k=\dfrac{x\sqrt{37}}{5}\sin(120^\circ-\theta).$ We can do a similar process to find $FC$, and adding them together gives (since $AF+FC=AC=\sqrt{37}$):
 \[\dfrac{x\sqrt{37}}{5}\sin(120^\circ-\theta)+\dfrac{x\sqrt{37}}{2\sqrt{3}}\sin(\theta+30^\circ)=\sqrt{37}.\]We can divide by $x\sqrt{37}$ to get
 \[\dfrac{1}{5}\sin(120^\circ-\theta)+\dfrac{1}{2\sqrt{3}}\sin(\theta+30^\circ)=\dfrac{1}{x}.\]We can expand everything to find that
 \[\dfrac{\sqrt{3}}{10}\cos\theta+\dfrac{1}{10}\sin\theta+\dfrac{1}{4}\sin\theta+\dfrac{\sqrt{3}}{12}\cos\theta=\dfrac{1}{x}.\]Combining terms:
 \[\dfrac{11\sqrt{3}}{60}\cos\theta+\dfrac{7}{20}\sin\theta=\dfrac{1}{x}.\]We are trying to find the minimum value of $x$, so we can try to find that maximum value of $\dfrac{1}{x}.$ To do this, we can take the derivative and set it equal to $0$. Doing this, we find that $\tan{x}=\dfrac{7\sqrt{3}}{11}.$ We can use this to find that $\cos\theta=\dfrac{11}{2\sqrt{67}}$ and $\sin\theta=\dfrac{7\sqrt{3}}{2\sqrt{67}}$.

 Substituting this back in gives
 \[\dfrac{1}{60}\left(\dfrac{121\sqrt{3}}{2\sqrt{67}}+\dfrac{147\sqrt{3}}{2\sqrt{67}}\right)=\dfrac{1}{x}.\]We can simplify to find that $\dfrac{1}{30}\sqrt{201}=\dfrac{1}{x}$, so $x=\dfrac{30}{\sqrt{201}}.$ We can simply use the formula for the area of an equilateral triangle to get
 \[\dfrac{\dfrac{900}{201}\sqrt{3}}{4}=\dfrac{75\sqrt{3}}{67}\rightarrow \boxed{145}.\]




\textbf{Solution 16}

Label the triangle as $\triangle ABC$ where $AB = 2\sqrt{3}, BC = 5, CA = \sqrt{37}$. Label the point where the equilateral triangle intersects with $BC$ as $D$ and the point where it intersects with $AC$ at $E$ and the last point as $F$. 

 Define points $G$ and $H$ on $EF$ and $BC$ respectively such that $\angle DGF = \angle EHD = 90^{\circ}$. 

 Let $\angle FDB = \theta$ and let $FG = GE = x$. Then, $FD = 2x$ and $FB = 2x\sin(\theta)$ and $DB = 2x\cos(\theta)$. 

 By simple angle chasing, we find $\angle EDH = 120^{\circ} - \theta$ and $\angle AFE = \theta + 30^{\circ}$. Thus, we have

 $DC = BC - BD = 5 - 2x\cos(\theta), AF = AB - BF = 2\sqrt{3} - 2x\sin(\theta), EH = 2x\sin(120^{\circ} - \theta)$. 

 We will now find the area of $\triangle ABC$ using these smaller triangles. 

 Area of FBD: This equals $\frac{1}{2} \cdot 2x\sin(\theta) \cdot 2x\cos(\theta) = 2x^2 \sin(\theta) \cos(\theta)$.

 Area of DEF: This equals $\frac{\sqrt{3}}{4} \cdot (2x)^2 = x^2 \sqrt{3}$.

 Area of DEC: This equals $\frac{1}{2} \cdot 2x\sin(120^{\circ} - \theta) \cdot (5 - 2x\cos(\theta)) = (5 - 2x\cos(\theta))(x\sin(120^{\circ} - \theta)$.

 Area of AFE: Use $\angle AFE$ as the central angle. This would give $\frac{1}{2} \cdot 2x \cdot (2\sqrt{3} - 2x\sin(\theta)) \cdot \sin(\theta + 30^{\circ}) = x(2\sqrt{3} - 2x\sin(\theta))(\sin(\theta + 30^{\circ}))$.

 Thus, the total area gives $2x^2 \sin(\theta) \cos(\theta) + x^2 \sqrt{3} + (5 - 2x\cos(\theta))(x\sin(120^{\circ} - \theta)) + x(2\sqrt{3} - 2x\sin(\theta))(\sin(\theta + 30^{\circ}))$.

 I am going to simplify the $x^2$ terms first. Factor out an $x^2$ from all the $x^2$ terms.

 This simplifies to

 $x^2(2\sin(\theta)\cos(\theta) + \sqrt{3} - 2\cos(\theta)\sin(120^{\circ} - \theta) - 2\sin(\theta)\sin(\theta + 30^{\circ}) = x^2(2\sin(\theta)\cos(\theta) + \sqrt{3} - 2\cos(\theta)(\frac{\sqrt{3}}{2} \cos(\theta) + \frac{1}{2} \sin(\theta)) - 2\sin(\theta)(\frac{\sqrt{3}}{2} \sin(\theta) + \frac{1}{2} \cos(\theta))$.

 This gives

 $x^2(2\sin(\theta)\cos(\theta) + \sqrt{3} - \cos^2(\theta) \sqrt{3} - \cos(\theta)\sin(\theta) - \sin^2(\theta) \sqrt{3} - \sin(\theta)\cos(\theta)) = x^2(2\sin(\theta)\cos(\theta) + \sqrt{3} - \sqrt{3}(\sin^2(\theta) + \cos^2(\theta)) - 2\sin(\theta)\cos(\theta)) = x^2(0) = 0$. 

 Finally, we're just left with all the $x$ terms namely

 $5x\sin(120^{\circ} - \theta) + 2x\sqrt{3} \sin(\theta + 30^{\circ})$

 Recall that this entire expression is just the area of $\triangle ABC$ which equals $\frac{1}{2} \cdot AB \cdot BC = \frac{1}{2} \cdot 2\sqrt{3} \cdot 5 = 5\sqrt{3}$. Thus, we get

 $5x(\frac{\sqrt{3}}{2} \cos(\theta) + \frac{1}{2} \sin(\theta)) + 2x\sqrt{3}(\frac{\sqrt{3}}{2} \sin(\theta) + \frac{1}{2} \cos(\theta)) = 5\sqrt{3} \implies 11x\sin(\theta) + 7x\sqrt{3}\cos(\theta) = 10\sqrt{3}$.

 Therefore

 $x = \frac{10\sqrt{3}}{11\sin(\theta) + 7\sqrt{3} \cos(\theta)}$

 Recall that the area of the equilateral triangle is $x^2 \sqrt{3}$ and we want to minimize this. Thus, we want to maximize the denominator of $x$. By Cauchy-Schwarz Inequality

 $(11^2 + (7\sqrt{3})^2)(\sin^2(\theta) + \cos^2(\theta)) \geq (11\sin(\theta) + 7\sqrt{3} \cos(\theta))^2 \implies (11\sin(\theta) + 7\sqrt{3} \cos(\theta)) \leq 2\sqrt{67}$. 

 Thus, $x \geq \frac{10\sqrt{3}}{2\sqrt{67}} = \frac{5\sqrt{3}}{\sqrt{67}} \implies x^2 \geq \frac{75}{67}$.

 Thus the minimum area is $x^2 \sqrt{3} = \frac{75\sqrt{3}}{67} \implies 75 + 3 + 67 = \boxed{145}$



\end{enumerate}